% ============================================================================
% NASKAH DRAF v1 (27 Oktober 2026)
% Knowledge-Based System Adaptif untuk Perhitungan PPh 21 dengan Akuisisi
% Pengetahuan Berbantuan LLM dari Regulasi dan Kebijakan Perusahaan
%
% Disusun dari paper/draft/skeleton.tex (versi origin/main b2aec05) dan
% feedback_and_revision/revision_27_10_26.md. Skeleton tidak diubah; berkas
% ini adalah salinan kerja terpisah.
%
% Format : IEEEtran (journal), seperti skeleton
% Bahasa : Indonesia; istilah teknis asing ditulis miring lewat \en{}
% Kompilasi: pdflatex, bibtex, lalu pdflatex dua kali
%            (atau: tectonic -X compile main.tex)
%
% Penanda draf (matikan semuanya dengan \rancanganfalse):
%   \status{}      kotak ungu  : status implementasi atau data yang sebenarnya
%   \terbuka{}     kotak oranye: keputusan yang belum diambil peneliti
%   \definisikan{} kotak hijau : hal yang masih harus didefinisikan
%   \catatan{}     kotak abu   : catatan penyusunan draf
%   \isian{}       teks abu    : angka atau teks yang menunggu hasil
%   \rujuk{}       teks merah  : rujukan yang belum ada di bibliografi
% ============================================================================
\documentclass[journal,a4paper]{IEEEtran}
\usepackage[english,indonesian]{babel}
\usepackage{cite}
\usepackage{amsmath,amssymb,amsfonts,amsthm,mathtools}
\usepackage{graphicx}
\usepackage{booktabs}
\usepackage{tabularx}
\usepackage{array}
\usepackage{multirow}
\usepackage{enumitem}
\usepackage{algorithm}
\usepackage{algpseudocode}
\usepackage{xcolor}
\usepackage{framed}
\usepackage{listings}
\usepackage{tikz}
\usetikzlibrary{positioning,arrows.meta,fit,calc,shapes.geometric}
\usepackage{url}
\usepackage[hidelinks]{hyperref}

\newcolumntype{L}{>{\raggedright\arraybackslash}X}
\newcolumntype{P}[1]{>{\raggedright\arraybackslash}p{#1}}

% Istilah asing (tidak diterjemahkan) ditulis miring.
\newcommand{\en}[1]{\textit{#1}}

% Penamaan dalam Bahasa Indonesia
\renewcommand{\IEEEkeywordsname}{Kata Kunci}
\addto\captionsindonesian{\renewcommand{\refname}{Daftar Pustaka}\renewcommand{\abstractname}{Abstrak}%
  \renewcommand{\figurename}{Gambar}\renewcommand{\tablename}{Tabel}}
\floatname{algorithm}{Algoritma}
\algrenewcommand\algorithmicrequire{\textbf{Masukan:}}
\algrenewcommand\algorithmicensure{\textbf{Keluaran:}}
\algrenewcommand\algorithmicforall{\textbf{untuk setiap}}
\algrenewcommand\algorithmicfor{\textbf{untuk}}
\algrenewcommand\algorithmicif{\textbf{jika}}
\algrenewcommand\algorithmicthen{\textbf{maka}}
\algrenewcommand\algorithmicelse{\textbf{selain itu}}
\algrenewcommand\algorithmicend{\textbf{akhir}}
\algrenewcommand\algorithmicdo{\textbf{lakukan}}
\algrenewcommand\algorithmicreturn{\textbf{kembalikan}}
\newtheorem{definisi}{Definisi}

% Notasi: x = kasus uji, t = masa pajak, \omega = tanggal pengetahuan,
% m = jumlah kandidat, c = jumlah kandidat benar, n = indeks peristiwa.
\newcommand{\K}{\mathcal{K}}
\newcommand{\hasil}{\mathit{hasil}}
\newcommand{\emas}{\star}
\newcommand{\SG}{S^{G}}
\newcommand{\SE}{S^{E}}

% Listing Turtle/SHACL/SPARQL
\lstdefinelanguage{rdf}{
  morekeywords={PREFIX,SELECT,DISTINCT,WHERE,FILTER,a},
  sensitive=true,
  comment=[l]{\#},
  morestring=[b]"
}
\lstset{
  basicstyle=\ttfamily\scriptsize,
  keywordstyle=\color{blue!60!black}\bfseries,
  stringstyle=\color{green!40!black},
  commentstyle=\color{gray},
  frame=single, framerule=0.3pt, rulecolor=\color{gray!50},
  breaklines=true, columns=fullflexible, xleftmargin=2pt, xrightmargin=2pt,
  aboveskip=4pt, belowskip=2pt
}

% Warna lapisan (sama dengan slide progres)
\definecolor{lapreg}{HTML}{5B8DEF}
\definecolor{lapprs}{HTML}{F2A15A}

% ---------------------------------------------------------------------------
% Penanda draf
% ---------------------------------------------------------------------------
\newif\ifrancangan \rancangantrue
\newcommand{\kotak}[3]{\ifrancangan\par\begingroup\scriptsize
  \def\FrameCommand{\fboxsep=3pt\fboxrule=0.4pt\fcolorbox{#1!60!black}{#1!6}}%
  \MakeFramed{\advance\hsize-\width\FrameRestore}\noindent\textbf{#2}\par #3\endMakeFramed\endgroup\fi}
\newcommand{\status}[1]{\kotak{violet}{Status}{#1}}
\newcommand{\terbuka}[1]{\kotak{orange}{Keputusan terbuka}{#1}}
\newcommand{\definisikan}[1]{\kotak{teal}{Harus didefinisikan}{#1}}
\newcommand{\catatan}[1]{\kotak{gray}{Catatan draf}{#1}}
\newcommand{\isian}[1]{\ifrancangan\textcolor{gray}{[#1]}\else[#1]\fi}
\newcommand{\rujuk}[1]{\textcolor{red}{[rujukan: #1]}}
\newcommand{\gambar}[1]{\fbox{\parbox[c][2.6cm][c]{0.94\columnwidth}{\centering\scriptsize #1}}}
\newcommand{\gambarlebar}[1]{\fbox{\parbox[c][3.2cm][c]{0.96\textwidth}{\centering\scriptsize #1}}}

\begin{document}
% Lebih dari enam penulis ditulis "et al." (gaya IEEE); entri kontrol ada di rujukan_tambahan.bib
\bstctlcite{IEEEexample:BSTcontrol}
\title{\en{Knowledge-Based System} Adaptif untuk Perhitungan PPh 21 dengan Akuisisi Pengetahuan Berbantuan LLM dari Regulasi dan Kebijakan Perusahaan}
\author{[Nama Penulis]%
\thanks{Draf v1, 27 Oktober 2026. Bagian I--VI berupa prosa; abstrak dan Bagian VII--IX menunggu hasil eksperimen.}}
\maketitle

% ============================================================================
% ABSTRAK (DRAF): kalimat hasil dan implikasi menunggu E1--E4.
% Batas 150--250 kata, satu paragraf, tanpa rumus, tabel, atau rujukan.
% ============================================================================
\begin{abstract}
Perhitungan PPh 21 harus mengikuti perubahan regulasi dan kebijakan perusahaan, tetapi aturannya umumnya tertanam di kode program atau \en{spreadsheet}. Akibatnya setiap perubahan menuntut pemrograman ulang dan dapat menggeser hasil periode lama tanpa terdeteksi. Penelitian formalisasi hukum pajak berbantuan \en{large language model} (LLM) membangun basis aturan dari nol untuk keadaan hukum pada satu waktu, sehingga belum diketahui apakah pembaruan berulang oleh LLM menjaga hasil periode lama dan apakah kesalahannya merambat dari satu perubahan ke perubahan berikutnya. Penelitian ini mengusulkan \en{knowledge-based system} adaptif yang menyimpan pengetahuan sebagai \en{legal knowledge graph} temporal dua lapis, yaitu lapisan regulasi dan lapisan kebijakan perusahaan, dan menjalankan seluruh perhitungannya dengan \en{inference engine} deterministik. Pembaruan \en{knowledge base} disusun LLM dari dokumen regulasi atau kebijakan, lalu disaring lima gerbang verifikasi dan dipilih lewat konsensus perilaku. Evaluasinya memutar ulang perubahan regulasi PPh 21 pegawai tetap yang nyata pada 2023--2026, termasuk pergantian metode ke tarif efektif rata-rata, menguji perubahan yang belum pernah dilihat perancang, dan membandingkan usulan dengan pembanding serta ablasi. \isian{Kalimat hasil: RA pada perubahan tak terlihat, SER dan CAR beserta selang kepercayaan 95\%, OAR dan kesalahan terpicu pada \en{regulation replay}, serta \en{recall} analisis dampak.} \isian{Kalimat implikasi bagi sistem payroll dan formalisasi peraturan yang sering berubah.}
\end{abstract}

\begin{IEEEkeywords}
Akuisisi pengetahuan, \en{knowledge-based system}, \en{knowledge graph}, \en{large language model}, perubahan regulasi, PPh Pasal 21, \en{rules as code}
\end{IEEEkeywords}

\catatan{Draf ini mengikuti urutan \texttt{paper/draft/skeleton.tex}. Bagian I--III ditulis lengkap. Bagian IV--VI sudah berupa prosa, tetapi memuat keputusan yang belum diambil (kotak oranye) dan butir yang belum didefinisikan (kotak hijau). Abstrak dan Bagian VII--IX hanya kerangka, karena eksperimen belum dijalankan. Kotak ungu menyatakan status yang sebenarnya, yaitu bagian mana yang baru rancangan dan bagian mana yang sudah ada di prototipe. Semua penanda hilang dengan \texttt{\textbackslash rancanganfalse}.}

% ============================================================================
\section{Pendahuluan}
\label{sec:pendahuluan}
% ============================================================================

\IEEEPARstart{P}{Ph} Pasal 21 dipotong pemberi kerja setiap bulan dari penghasilan pegawai, sehingga kesalahan perhitungannya langsung muncul di setiap slip gaji. Regulasinya juga sering berubah. Sampai tahun pajak 2023, pemotongan bulanan memakai metode disetahunkan menurut PER-16/PJ/2016~\cite{per16_2016}. Sejak 1 Januari 2024, PP 58/2023~\cite{pp58_2023} dan PMK 168/2023~\cite{pmk168_2023} menggantinya dengan tarif efektif rata-rata (TER) bulanan dan perhitungan ulang setahun di masa pajak terakhir. Pada Februari 2025 terbit fasilitas PPh 21 ditanggung pemerintah (DTP) untuk industri padat karya yang berlaku surut sejak Januari 2025~\cite{pmk10_2025}. Fasilitas ini diperluas ke sektor pariwisata pada Oktober 2025~\cite{pmk72_2025} dan diperpanjang untuk 2026~\cite{pmk105_2025}. Batas upah iuran Jaminan Pensiun, yang menentukan besarnya pengurang penghasilan, juga naik setiap Maret. Kandidat korpus penelitian ini mencatat sembilan peristiwa perubahan yang memengaruhi PPh 21 pegawai tetap pada 2023--2026, atau dua sampai tiga peristiwa per tahun, dengan jenis yang beragam, yaitu perubahan parameter, klasifikasi penghasilan, fasilitas baru, dan pergantian metode perhitungan (Tabel~\ref{tab:korpus}).

\status{Angka sembilan peristiwa berasal dari daftar kandidat di catatan revisi 27 Oktober (empat kenaikan batas upah JP, PMK 66/2023, pergantian ke TER, PMK 10/2025, PMK 72/2025, PMK 105/2025). Daftar ini belum diverifikasi sebagai korpus final. Bukti biaya perubahan pada praktik saat ini (pembanding B1) akan ditambahkan setelah Tahap 10.}

Dalam praktik, aturan pajak dan kebijakan payroll perusahaan tertanam bercampur di kode program atau rumus \en{spreadsheet}. Audit pendahuluan penelitian ini atas \en{spreadsheet} payroll sebuah perusahaan, yang selanjutnya disebut Perusahaan X, menemukan konstanta regulasi tertanam di 254 sel rumus. \en{Spreadsheet} itu juga mencatat lembur sebagai penghasilan tidak teratur, padahal PER-16/PJ/2016 dan PMK 168/2023 menggolongkannya sebagai penghasilan teratur. Survei pendahuluan atas tujuh pustaka dan aplikasi PPh 21 \en{open source} menemukan kesalahan terverifikasi di ketujuhnya. Contohnya lapisan tarif Pasal 17 yang diperlakukan sebagai lebar lapisan, bukan batas kumulatif, dan lebih bayar di masa pajak terakhir yang dipotong menjadi nol. Pada salah satu aplikasi, perhitungan TER baru ditambahkan sekitar 30 bulan setelah TER berlaku, sebagai halaman terpisah di samping kode metode lama. Pola seperti ini menimbulkan tiga masalah. Pertama, setiap perubahan regulasi menuntut pemrograman ulang. Kedua, tidak ada bukti bahwa perubahan itu tidak menggeser hasil periode lama, padahal hasil masa lampau tetap dibutuhkan untuk pembetulan dan pemeriksaan pajak. Ketiga, kebijakan perusahaan yang menjadi melanggar setelah regulasi berubah tidak terdeteksi.

\en{Knowledge-based system} (KBS) memisahkan pengetahuan dari kode, sehingga perubahan aturan cukup dilakukan pada \en{knowledge base}. Gerakan \en{rules as code} menerapkan gagasan yang sama pada peraturan perundang-undangan~\cite{merigoux2021catala,mowbray2023scaling,kowalski2022logicalenglish}. Pemisahan ini memindahkan beban ke pemeliharaan \en{knowledge base}. Setiap perubahan hukum harus dibaca, ditafsirkan, lalu ditulis ulang sebagai aturan oleh orang yang memahami hukum sekaligus representasinya. Masalah ini dikenal sebagai \en{knowledge acquisition bottleneck}. Metode akuisisi inkremental seperti \en{Ripple-Down Rules} menekan biaya setiap perubahan dengan menjaga kasus lama tetap benar, tetapi pakar manusia tetap harus menulis setiap perubahan~\cite{compton1990philosophical,compton2021rippledown}.

\en{Large language model} (LLM) dapat membantu tahap akuisisi ini, karena LLM mampu mengubah teks peraturan menjadi representasi terstruktur~\cite{janatian2023texttostructure}. Namun, LLM tidak andal untuk penalaran pajak yang menuntut ketepatan~\cite{blairstanek2023gpt3,nay2024taxattorneys}, dapat mengarang fakta hukum~\cite{dahl2024legalfictions,magesh2025hallucination}, dan sering menerapkan ketentuan yang sudah diganti amandemen~\cite{prior2026asking,huang2026wrong}. Karena itu, penelitian formalisasi pajak berbantuan LLM memakai LLM untuk menyusun aturan, lalu menyerahkan perhitungan kepada \en{engine} deterministik~\cite{veeramani2026neurosymbolic,yadamsuren2026formalization,lorenzo2025translating,han2026autoformalization,jurayj2026tax}. Karya-karya itu membangun basis aturan dari nol untuk keadaan hukum pada satu waktu. Gagasan memakai LLM untuk memperbarui perangkat lunak pajak versi sebelumnya dari publikasi otoritas pajak sudah diajukan, tetapi baru sebagai kajian tantangan teknis~\cite{goganikhiabani2025maintaining}. Di sisi lain, penyuntingan pengetahuan LLM yang dilakukan berurutan terbukti dapat merusak pengetahuan yang sudah ada~\cite{gupta2024forgetting,gu2024harms}. Pada konsolidasi undang-undang dengan LLM, kemiripan teks hasil konsolidasi mencapai 93--99\%, tetapi kecocokan persisnya hanya 50,3\% untuk amandemen satu langkah dan 20,51\% untuk rangkaian amandemen~\cite{prior2025risks}. Kemiripan yang tinggi menyembunyikan hasil yang salah, dan kesalahan bertambah ketika perubahan berantai. Sejauh penelusuran penulis, belum ada penelitian yang mengukur tiga hal pada basis aturan yang dieksekusi, yaitu apakah pembaruan berulang dari LLM menjaga hasil periode lama, apakah kesalahan merambat dari satu perubahan ke perubahan berikutnya, dan bagaimana dampak perubahan regulasi terhadap lapisan kebijakan perusahaan.

Penelitian ini mengusulkan KBS adaptif untuk perhitungan PPh 21 pegawai tetap (Gambar~\ref{fig:ikhtisar}). Pengetahuan disimpan sebagai \en{legal knowledge graph} temporal dua lapis yang memisahkan regulasi pemerintah dari kebijakan perusahaan. Seluruh aritmetika dijalankan \en{inference engine} deterministik yang mencatat jejak setiap angka sampai pasal sumbernya. Ketika dokumen regulasi atau kebijakan baru terbit, LLM menyusun beberapa kandidat pembaruan \en{knowledge base}, lima gerbang verifikasi menyaringnya, dan sistem memilih satu kandidat lewat konsensus perilaku atau menolak memilih (\en{abstain}). Kebenaran sistem diukur dari hasil perhitungan sampai rupiah dengan memutar ulang seluruh perubahan regulasi yang nyata secara berurutan, yang disebut \en{regulation replay}.

\begin{figure*}[t]
\centering
\begin{tikzpicture}[
  font=\scriptsize,
  node distance=5mm and 4.2mm,
  langkah/.style={draw=black!70, rounded corners=2pt, align=center, text width=22mm, minimum height=10mm, fill=white},
  llm/.style={langkah, fill=violet!7},
  gerbang/.style={langkah, fill=teal!8},
  manusia/.style={langkah, fill=gray!10},
  lapis/.style={draw=black!60, rounded corners=1.5pt, align=center, text width=30mm, minimum height=6mm},
  panah/.style={-{Stealth[length=1.8mm]}, semithick},
  ket/.style={font=\tiny, align=center, fill=white, inner sep=1pt}
]
% Baris atas: alur akuisisi
\node[manusia] (dok) {Dokumen regulasi\\atau kebijakan baru};
\node[llm, right=of dok] (rel) {Penilaian relevansi\\berkutipan (LLM)};
\node[llm, right=of rel] (kand) {$m$ kandidat\\pembaruan $\Delta$ (LLM)};
\node[gerbang, right=of kand] (grb) {Gerbang\\G1--G5};
\node[gerbang, right=of grb] (pilih) {Konsensus perilaku\\atau \en{abstain}};
\node[manusia, right=of pilih] (tinjau) {Peninjau manusia\\dan laporan dampak};
% Baris bawah: knowledge graph dan perhitungan
\node[lapis, fill=lapreg!25] (reg) at ($(kand.south)!0.5!(grb.south)+(0,-17mm)$) {Lapisan regulasi};
\node[lapis, fill=lapprs!30, below=1.5mm of reg] (prs) {Lapisan perusahaan};
\node[draw=black!70, rounded corners=2pt, fit=(reg)(prs), inner sep=2.5mm,
      label={[font=\scriptsize\itshape, align=center]below:Knowledge graph temporal\\(RDF/OWL, SHACL, SPARQL)}] (kg) {};
\node[langkah, fill=blue!6] (eng) at (pilih |- kg) {\en{Inference engine}\\deterministik};
\node[langkah, fill=green!6] (out) at (tinjau |- kg) {PPh 21 per masa dan\\jejak sampai pasal};
\node[manusia, below=4mm of eng, minimum height=7mm] (data) {Data pegawai, $t$, $\omega$};
% Panah alur
\draw[panah] (dok) -- (rel);
\draw[panah] (rel) -- (kand);
\draw[panah] (kand) -- (grb);
\draw[panah] (grb) -- (pilih);
\draw[panah] (pilih) -- (tinjau);
\draw[panah] (grb.north) to[bend right=40] node[ket, above]{umpan balik perbaikan} (kand.north);
\draw[panah, dashed] (rel.south) |- node[ket, pos=0.3]{tidak relevan:\\$\Delta$ kosong} (kg.west);
\draw[panah] (kg.north -| kand) -- node[ket]{konteks} (kand.south);
\draw[panah] (tinjau.south) -- ++(0,-4mm) -| node[ket, pos=0.3]{pembaruan diterapkan sebagai versi baru} ($(kg.north)+(9mm,0)$);
\draw[panah] (kg.east) -- (eng.west);
\draw[panah] (eng) -- (out);
\draw[panah] (data) -- (eng);
\end{tikzpicture}
\caption{Ikhtisar sistem yang diusulkan. Baris atas adalah alur akuisisi pengetahuan. Baris bawah adalah \en{knowledge base} dua lapis dan \en{inference engine} yang menghitung PPh 21. LLM hanya mengusulkan pembaruan dan tidak pernah menghitung pajak.}
\label{fig:ikhtisar}
\end{figure*}

Penelitian ini menjawab empat pertanyaan penelitian (RQ).
\begin{enumerate}[label=\textbf{RQ\arabic*},leftmargin=*,nosep]
  \item \textbf{Adaptivitas representasi.} Sejauh mana representasi pengetahuan temporal dua lapis dapat menyatakan perubahan regulasi dan kebijakan, termasuk jenis perubahan yang tidak diketahui saat representasi dirancang, tanpa mengubah kode \en{engine}?
  \item \textbf{Akuisisi satu langkah.} Seberapa sering akuisisi berbantuan LLM menerapkan pembaruan yang benar, seberapa sering ia menerapkan pembaruan salah tanpa terdeteksi, dan berapa besar peran gerbang verifikasi pada kedua angka itu?
  \item \textbf{Adaptivitas operasional berantai.} Pada rangkaian perubahan regulasi yang nyata, berapa peristiwa yang berhasil diterapkan otomatis, dan apakah kesalahan merambat dari satu peristiwa ke peristiwa berikutnya?
  \item \textbf{Dampak lintas lapis.} Seberapa lengkap dan tepat sistem mengidentifikasi kebijakan perusahaan yang terdampak atau menjadi melanggar setelah regulasi berubah?
\end{enumerate}

Kontribusi penelitian ini adalah sebagai berikut.
\begin{enumerate}[label=\textbf{K\arabic*},leftmargin=*,nosep]
  \item Representasi pengetahuan temporal dua lapis untuk perhitungan pajak penghasilan pegawai, dengan resolusi konflik antarlapis dan analisis dampak lintas lapis (RQ1, RQ4).
  \item Metode akuisisi pengetahuan inkremental berbantuan LLM dengan gerbang verifikasi, gerbang non-vakuitas, dan pemilihan kandidat berbasis konsensus perilaku (RQ2).
  \item Definisi adaptivitas representasional dan operasional, serta protokol \en{regulation replay} yang memisahkan kesalahan bawaan, kesalahan baru, dan kesalahan yang dipicu perambatan, beserta bukti empiris pada rangkaian perubahan PPh 21 yang nyata (RQ3).
  \item Artefak terbuka berupa \en{knowledge base}, korpus perubahan beserta pembaruan emas, \en{test suite} dengan \en{ground truth}, dan kode evaluasi.
\end{enumerate}
Tabel~\ref{tab:lacak} di Bagian~\ref{sec:evaluasi} menghubungkan setiap kontribusi dengan RQ, hipotesis, dan metrik utamanya. Penelitian ini tidak mengklaim pola ``LLM menyusun aturan, \en{engine} deterministik menghitung'' sebagai kebaruan, karena pola itu sudah dipakai karya terdahulu (Bagian~\ref{sec:terkait}). Gagasan pemeliharaan perangkat lunak pajak dengan LLM juga sudah diajukan. Kebaruannya terletak pada pembaruan inkremental \en{knowledge base} yang dapat dieksekusi dan dijaga gerbang verifikasi, yang dievaluasi pada rangkaian amandemen nyata, beserta pengukuran perambatan kesalahan dan dampaknya terhadap lapisan kebijakan perusahaan.

Sisa naskah disusun sebagai berikut. Bagian~\ref{sec:masalah} menjelaskan perhitungan PPh 21 pegawai tetap, batas kewenangan kebijakan perusahaan, jenis perubahan regulasi, dan rumusan masalah. Bagian~\ref{sec:terkait} membahas penelitian terkait. Bagian~\ref{sec:representasi} menetapkan representasi pengetahuan dan \en{inference engine}. Bagian~\ref{sec:akuisisi} mendefinisikan pembaruan dan adaptivitas, lalu menjelaskan alur akuisisi berbantuan LLM. Bagian~\ref{sec:evaluasi} menguraikan metodologi evaluasi. Bagian~\ref{sec:hasil} dan~\ref{sec:pembahasan} melaporkan dan membahas hasil, dan Bagian~\ref{sec:kesimpulan} menutup naskah.

% ============================================================================
\section{Latar belakang dan rumusan masalah}
\label{sec:masalah}
% ============================================================================

Bagian ini memberi pembaca dari luar Indonesia pengetahuan domain secukupnya untuk memahami mengapa perubahan PPh 21 sulit diikuti, lalu merumuskan masalahnya.

\subsection{Perhitungan PPh 21 pegawai tetap}
\label{subsec:pph21}

PPh 21 pegawai tetap dihitung dari penghasilan bruto. Bruto terdiri atas penghasilan teratur, seperti gaji, tunjangan, dan uang lembur, dan penghasilan tidak teratur, seperti tunjangan hari raya (THR) dan bonus. Premi jaminan kecelakaan kerja (JKK), jaminan kematian (JKM), dan jaminan kesehatan yang dibayar pemberi kerja ikut menjadi bruto, sedangkan iuran jaminan hari tua (JHT) dan jaminan pensiun (JP) yang dibayar pemberi kerja tidak~\cite{pmk168_2023}. Penghasilan neto adalah bruto dikurangi biaya jabatan, yaitu 5\% dari bruto dengan batas Rp500.000 sebulan dan Rp6.000.000 setahun, serta iuran JHT dan JP yang dibayar pegawai. Iuran JP pegawai sebesar 1\% dihitung dari upah paling tinggi sebesar batas upah yang ditetapkan setiap tahun. Penghasilan kena pajak (PKP) adalah neto setahun dikurangi penghasilan tidak kena pajak (PTKP), yaitu Rp54 juta untuk diri wajib pajak, ditambah Rp4,5 juta bila kawin dan Rp4,5 juta untuk setiap tanggungan sampai tiga orang~\cite{pmk101_2016}. PKP dibulatkan ke bawah hingga ribuan rupiah. PPh setahun dihitung dengan tarif progresif Pasal 17 UU PPh, yang sejak UU HPP terdiri atas lima lapisan, yaitu 5\% sampai Rp60 juta, 15\% sampai Rp250 juta, 25\% sampai Rp500 juta, 30\% sampai Rp5 miliar, dan 35\% di atasnya~\cite{uu7_2021}.

Cara memotong PPh 21 setiap bulan berubah pada 2024. Sampai tahun pajak 2023, pemberi kerja menyetahunkan penghasilan teratur setiap bulan, menghitung PPh setahun, lalu membaginya per bulan. Penghasilan tidak teratur dipajaki terpisah sebagai selisih PPh setahun dengan dan tanpa penghasilan itu~\cite{per16_2016}. Sejak 2024, potongan bulan Januari sampai November adalah TER dikali bruto bulan itu. TER terbagi menjadi kategori A, B, dan C menurut status PTKP, dengan 44, 40, dan 41 lapisan tarif bulanan~\cite{pp58_2023}. Pada masa pajak terakhir, yaitu Desember atau bulan pegawai berhenti bekerja, pemberi kerja menghitung PPh setahun dengan tarif Pasal 17, lalu memotong selisihnya dengan jumlah yang sudah dipotong~\cite{pmk168_2023}. Selisih ini dapat negatif. Dalam hal itu, kelebihan potongan dikembalikan kepada pegawai.

Pemberi kerja dapat memotong PPh 21 dari penghasilan pegawai, memberi tunjangan pajak sebesar PPh 21 itu sendiri (\en{gross-up}), atau menanggungnya. Pada \en{gross-up}, tunjangan pajak adalah PPh atas bruto yang memuat tunjangan itu, sehingga nilainya didefinisikan secara melingkar. Pada 2025 dan 2026, pemerintah juga menanggung PPh 21 pegawai di sektor tertentu yang bruto tetap dan teraturnya paling banyak Rp10 juta sebulan. PPh yang ditanggung pemerintah dibayarkan tunai kepada pegawai dan tidak dihitung sebagai penghasilan~\cite{pmk10_2025,pmk72_2025,pmk105_2025}. Tabel~\ref{tab:ruang-pajak} merangkum dimensi yang menentukan besarnya pajak.

\begin{table}[t]
\centering
\caption{Dimensi perhitungan PPh 21 pegawai tetap yang dicakup}
\label{tab:ruang-pajak}
\scriptsize
\begin{tabularx}{\columnwidth}{@{}P{2.0cm} L@{}}
\toprule
\textbf{Dimensi} & \textbf{Nilai} \\
\midrule
Rezim & Disetahunkan (sampai 2023); TER (sejak 2024); TER dengan fasilitas DTP (2025--2026) \\
Lapisan Pasal 17 & 5\%, 15\%, 25\%, 30\%, 35\% (empat lapisan sebelum UU HPP, hanya untuk data pengembangan) \\
Kategori TER & A, B, C dengan 44, 40, dan 41 lapisan tarif bulanan \\
Status PTKP & TK/0--TK/3, K/0--K/3; menentukan kategori TER \\
Metode & Dipotong (\en{gross}), \en{gross-up}, ditanggung pemberi kerja \\
Pola kerja & Setahun penuh, masuk di tengah tahun, keluar di tengah tahun \\
Penghasilan & Teratur; tidak teratur (THR, bonus) \\
\bottomrule
\end{tabularx}
\end{table}

\status{Jumlah lapisan TER 44, 40, dan 41 berasal dari \texttt{dataset/01\_regulasi/tables/ter\_bulanan.csv}, yang sudah diekstraksi dua kali, dicek visual pada beberapa halaman lampiran PP 58/2023, dan cocok dengan contoh PMK 168/2023. Catatan revisi masih meminta cek ulang penuh ke lampiran PP 58/2023 (Tahap 2).}

Penelitian ini terbatas pada pegawai tetap yang menerima penghasilan dari satu pemberi kerja, pada tahun pajak 2023--2026. Pegawai tidak tetap, bukan pegawai, pensiunan, pegawai dengan lebih dari satu pemberi kerja, pindah cabang, fasilitas khusus Ibu Kota Nusantara, dan penilaian natura berada di luar lingkup. Nilai natura diterima sebagai masukan. Dimensi di Tabel~\ref{tab:ruang-pajak} saling membatasi. Status PTKP menentukan kategori TER. Metode disetahunkan hanya berlaku sampai 2023, TER hanya sejak 2024, dan DTP hanya pada 2025--2026 untuk sektor dan batas bruto tertentu. Pegawai yang masuk atau keluar di tengah tahun mengubah masa pajak terakhir dan cara menghitung PPh setahun. Ruang uji di Bagian~\ref{subsec:ruanguji} hanya membangkitkan kombinasi yang sah menurut batasan ini.

\terbuka{Prototipe saat ini menghitung \en{gross-up} hanya pada rezim TER. Apakah \en{gross-up} pada rezim disetahunkan (2023) masuk lingkup, atau dicatat sebagai batasan antardimensi?}

\subsection{Kebijakan perusahaan dan batas kewenangannya}
\label{subsec:kebijakan}

Perusahaan menetapkan sebagian besar komponen penghasilan pegawai lewat peraturan perusahaan, perjanjian kerja bersama (PKB), atau keputusan direksi. Contohnya cara memprorata gaji pegawai yang masuk di tengah bulan, cara menghitung THR, besaran tunjangan, dan dasar upah iuran BPJS. Kebijakan seperti ini menentukan besarnya bruto, sehingga ikut menentukan PPh 21. Perusahaan juga boleh memilih di antara pilihan yang disediakan regulasi, misalnya menanggung PPh 21 pegawai atau memberinya tunjangan pajak.

Perusahaan tidak boleh menyimpangi ketentuan wajib. Contohnya klasifikasi jenis penghasilan. PER-16/PJ/2016 Pasal 1 angka 15 dan PMK 168/2023 Pasal 5 ayat (3) huruf a menggolongkan uang lembur sebagai penghasilan teratur~\cite{per16_2016,pmk168_2023}. Kebijakan payroll Perusahaan X mencatatnya sebagai penghasilan tidak teratur. Pada metode disetahunkan, penghasilan teratur ikut disetahunkan sedangkan penghasilan tidak teratur tidak, sehingga perbedaan klasifikasi ini mengubah besarnya potongan bulanan. Bila kebijakan perusahaan bertentangan dengan ketentuan wajib, ketentuan regulasi yang berlaku (\en{lex superior}), dan konflik itu harus tercatat. Di \en{spreadsheet} Perusahaan X, konflik ini tersembunyi di dalam rumus.

\subsection{Jenis perubahan regulasi}
\label{subsec:jenis-perubahan}

Tabel~\ref{tab:jenis} memuat taksonomi jenis perubahan yang dipakai untuk menyusun korpus dan melaporkan hasil. Satu peristiwa dapat termasuk lebih dari satu jenis. Contohnya PMK 10/2025, yang menambah fasilitas baru sekaligus berlaku surut. Saat melaporkan tingkat keberhasilan, jenis yang mudah seperti perubahan parameter tidak digabung dengan jenis struktural seperti pergantian metode.

\begin{table}[t]
\centering
\caption{Taksonomi jenis perubahan}
\label{tab:jenis}
\scriptsize
\begin{tabularx}{\columnwidth}{@{}P{1.8cm} P{2.5cm} L@{}}
\toprule
\textbf{Jenis} & \textbf{Contoh} & \textbf{Wujud pada \en{knowledge base}} \\
\midrule
Parameter & Batas upah JP setiap Maret & Nilai parameter baru dengan masa berlaku baru \\
Tabel tarif & Tabel TER (PP 58/2023) & Versi tabel baru \\
Klasifikasi & Perlakuan natura dan kenikmatan sebagai penghasilan (PMK 66/2023) & Pemetaan jenis penghasilan ke kategori pajak berubah \\
Aturan atau fasilitas baru & DTP padat karya (PMK 10/2025) & Aturan baru, sering disertai isian data baru \\
Kondisi atau cakupan & Sektor pariwisata (PMK 72/2025) & Kondisi aturan yang ada diperluas atau dipersempit \\
Pergantian metode & Disetahunkan ke TER (2024) & Sekelompok aturan diganti, aturan lama diakhiri \\
Pencabutan & PMK 250/2008 dicabut PMK 168/2023 & Masa berlaku aturan diakhiri secara eksplisit \\
Berlaku surut & PMK 10/2025 (terbit Februari, berlaku Januari) & Tanggal mulai berlaku lebih awal dari tanggal terbit \\
Kebijakan perusahaan & Versi PKB baru & Aturan lapisan perusahaan berubah \\
\bottomrule
\end{tabularx}
\end{table}

\status{Taksonomi ini adalah usulan draf yang menambah satu jenis dari daftar minimum di \en{skeleton}, yaitu perubahan kondisi atau cakupan. Taksonomi ditetapkan final di Tahap 2, bersama taksonomi konflik dan taksonomi kesalahan LLM.}

\subsection{Rumusan masalah}
\label{subsec:rumusan}

Masukan masalah ini adalah \en{knowledge base} $\K_{n-1}$ yang benar untuk keadaan hukum sebelum peristiwa perubahan ke-$n$, beserta dokumen sumber peristiwa itu. Keluarannya adalah $\K_n$ yang benar untuk keadaan hukum sesudahnya, dengan tiga syarat. Kode \en{engine} tidak berubah. Hasil yang tidak seharusnya berubah tetap sama, termasuk seluruh hasil untuk masa sebelum perubahan berlaku. Kebijakan perusahaan yang menjadi bertentangan dengan ketentuan wajib dilaporkan. Masalah ini berulang, karena $\K_n$ menjadi masukan untuk peristiwa berikutnya. Kesalahan yang lolos pada satu peristiwa dapat terbawa, atau memicu kesalahan baru, pada peristiwa sesudahnya. Bagian~\ref{subsec:adaptif} menyatakan syarat-syarat ini secara formal.

% ============================================================================
\section{Penelitian terkait}
\label{sec:terkait}
% ============================================================================

Bagian ini menunjukkan bahwa celah yang dirumuskan di Bagian~\ref{sec:pendahuluan} memang ada. Setiap subbagian ditutup dengan apa yang belum dikerjakan kelompok penelitian itu dan bagaimana penelitian ini berbeda. Tabel~\ref{tab:posisi} merangkum posisinya.

\subsection{\en{Rules as code} dan formalisasi peraturan pajak}
\label{subsec:rac}

Gerakan \en{rules as code} menuliskan peraturan sebagai spesifikasi yang dapat dieksekusi, sehingga aturan terpisah dari kode aplikasi. Catala adalah bahasa pemrograman untuk menerjemahkan teks undang-undang secara sistematis, termasuk pola kasus umum dan pengecualian yang dimodelkan dengan \en{default logic}. Evaluasinya pada pajak penghasilan Amerika Serikat dan tunjangan keluarga Prancis menemukan \en{bug} pada implementasi resmi tunjangan tersebut~\cite{merigoux2021catala}. Huttner dan Merigoux menegaskan bahwa proses pengembangan perangkat lunak konvensional sulit mengikuti arus perubahan peraturan~\cite{huttner2022catala}. Mlang menggantikan \en{compiler} warisan mesin pajak penghasilan Prancis yang makin sulit dipelihara setiap kali undang-undang berubah~\cite{merigoux2021mlang}. Ekosistem Catala kemudian menambahkan formalisasi aritmetika tanggal~\cite{monat2024date} dan pembangkitan kasus uji dengan \en{concolic execution}~\cite{goutagny2025cutecat}. Di luar Catala, DataLex merepresentasikan peraturan dalam bahasa aturan semi-alami yang dijalankan \en{inference engine} beserta penjelasannya~\cite{mowbray2023scaling,mowbray2023xai}, dan Logical English menawarkan bahasa alami terkendali yang dapat dieksekusi sebagai program logika~\cite{kowalski2022logicalenglish}. Witt dkk. membedakan validasi teknis dari \en{legal alignment}, yaitu kesamaan makna kode dengan undang-undang~\cite{witt2024encoding}. Pembedaan ini dipakai di penelitian ini untuk menyatakan batas klaim, karena kecocokan dengan acuan belum tentu berarti kesesuaian hukum. OpenFisca menyimpan parameter pajak dan tunjangan sebagai berkas bertanggal, sehingga nilai parameter mengikuti masa berlakunya~\cite{openfisca2026rules}.

Pada payroll, Dumancic dkk. menerjemahkan legislasi pajak dan jaminan sosial ke representasi formal yang dikompilasi menjadi \en{constraint} untuk mencari kebijakan payroll yang optimal~\cite{dumancic2021payroll}. Makowski dan Tr\k{a}bka memakai \en{Decision Model and Notation} dengan dukungan AI generatif untuk aturan pajak di kantor akuntan Swiss~\cite{makowski2025dmn}. Di Indonesia, sistem PPh 21 umumnya berupa aplikasi prosedural atau \en{spreadsheet} yang rumusnya ditanam di kode~\cite{ramadhani2021grossup,yulistyawati2022payroll,kusuma2026pph21}, dan penerapan TER menimbulkan lebih bayar pada pegawai dengan penghasilan tidak teratur~\cite{sherra2025ter}.

Pada semua karya ini, pengetahuan ditulis manusia. Parameter bertanggal sudah ada, tetapi pembaruan \en{knowledge base} oleh mesin mengikuti perubahan hukum belum dievaluasi. Penelitian ini juga memisahkan lapisan regulasi dari lapisan kebijakan pemberi kerja secara eksplisit, yang tidak menjadi fokus karya-karya tersebut.

\subsection{Akuisisi pengetahuan inkremental}
\label{subsec:inkremental}

\en{Ripple-Down Rules} menambah aturan secara inkremental ketika pakar menemukan kasus yang salah. Setiap aturan baru disertai \en{cornerstone case}, yaitu kasus yang membenarkan aturan itu, sehingga perubahan berikutnya tidak boleh merusak kasus-kasus tersebut~\cite{compton1990philosophical,compton2021rippledown}. Di bidang hukum, perubahan sistem hukum seperti abrogasi dan anulasi telah dimodelkan dengan \en{defeasible logic}~\cite{governatori2010changing}. Moguillansky dkk. membedakan saat suatu ketentuan masuk ke sistem hukum, saat ketentuan itu berlaku, dan saat ia menimbulkan akibat hukum, termasuk akibat yang berlaku surut, lalu merumuskan operator revisi yang menjaga konsistensi berbasis waktu~\cite{moguillansky2024temporalised}. Olivieri dkk. menunjukkan bahwa merevisi teori \en{defeasible logic} dengan perubahan seminimal mungkin adalah masalah yang sulit secara komputasi~\cite{olivieri2025revising}.

Gerbang regresi G4 dalam penelitian ini setara dengan \en{cornerstone case}, dan pemisahan masa berlaku dari tanggal pencatatan mengikuti pembedaan waktu pada revisi hukum temporal. Bedanya, pada karya-karya tersebut pakar manusia tetap menulis setiap perubahan, sedangkan penelitian ini mengukur seberapa jauh penulisan itu dapat diusulkan LLM tanpa merusak hasil lama.

\subsection{LLM untuk formalisasi hukum dan pajak}
\label{subsec:llm-formal}

LLM tidak andal bila diminta menalar atau menjawab pertanyaan hukum secara langsung. GPT-3 membuat kesalahan yang jelas pada penalaran undang-undang pajak, bahkan pada undang-undang sintetis sederhana~\cite{blairstanek2023gpt3}. Akurasi LLM dalam menerapkan hukum pajak meningkat antargenerasi model tetapi belum setara ahli~\cite{nay2024taxattorneys}. LLM publik berhalusinasi pada 58\% sampai 88\% pertanyaan faktual yang dapat diverifikasi tentang putusan pengadilan federal Amerika Serikat~\cite{dahl2024legalfictions}, dan alat riset hukum komersial berbasis \en{retrieval-augmented generation} masih berhalusinasi pada 17\% sampai 33\% kasus~\cite{magesh2025hallucination}.

Karena itu, kelompok karya terdekat memakai LLM untuk menyusun representasi, bukan untuk menghitung. Pada studi Janatian dkk., sekitar 60\% struktur yang dihasilkan LLM dari teks hukum dinilai setara atau lebih baik daripada hasil manual~\cite{janatian2023texttostructure}. Karya 2025--2026 menerapkan pola ini pada pajak. NTOL menggabungkan LLM dengan penalar simbolik deterministik di atas \en{knowledge graph} statuta pajak penghasilan India. Pada 20 skenario, akurasinya 80\%, dibanding 75\% untuk LLM saja, dan sebagian besar kesalahannya berasal dari cakupan graf yang belum lengkap~\cite{veeramani2026neurosymbolic}. Yadamsuren dkk. memakai LLM untuk menerjemahkan satu pasal \en{Internal Revenue Code} Amerika Serikat ke Prolog, lalu mendeteksi inkonsistensi undang-undang secara deterministik~\cite{yadamsuren2026formalization}. Lorenzo dkk. menyusun \en{benchmark} dan kerangka evaluasi untuk penerjemahan hukum pajak ke Catala dengan LLM~\cite{lorenzo2025translating}. Graus membangkitkan model keputusan yang dapat dieksekusi dari peraturan lingkungan dan tata ruang Belanda. Masukan terstruktur menaikkan kemiripan hasilnya dengan model produksi, tetapi kesetaraan fungsionalnya pada skenario uji masih terbatas~\cite{graus2026executable}. Law\&Order mengubah formulir dan petunjuk pengisian pajak menjadi program simbolik lewat sintesis LLM, verifikasi per sel, dan perbaikan berulang. Pada 51 SPT \en{holdout}, ketepatannya 100\%, sedangkan LLM terbaik tanpa kerangka itu hanya 66\%~\cite{han2026autoformalization}. Jurayj dkk. memadukan LLM dengan program logika untuk penalaran pajak~\cite{jurayj2026tax}. Pada pengujian perangkat lunak pajak, LLM dipakai untuk menurunkan spesifikasi metamorfik dari dokumen pajak~\cite{srinivas2023fewshot,khiabani2026agentic}.

Gagasan yang paling dekat dengan penelitian ini diajukan Gogani-Khiabani dkk. Mereka mengusulkan LLM untuk mengekstraksi perubahan kode dari publikasi otoritas pajak Amerika Serikat dan menerapkannya pada versi perangkat lunak pajak sebelumnya. Abstraknya membahas tantangan teknis gagasan itu dan tidak melaporkan evaluasi atas rangkaian perubahan~\cite{goganikhiabani2025maintaining}.

Karya-karya lain di kelompok ini memformalkan hukum untuk keadaan hukum pada satu waktu. Sejauh dapat dinilai dari abstraknya, belum ada yang mengevaluasi pembaruan berulang atas basis aturan yang sudah berjalan, perambatan kesalahan antarpembaruan, atau lapisan kebijakan pemberi kerja. Ketiga hal itu menjadi fokus penelitian ini.

\subsection{Pembaruan berulang oleh LLM dan evaluasinya}
\label{subsec:pembaruan-llm}

Penelitian penyuntingan pengetahuan LLM mengubah fakta yang tersimpan di bobot model. Suntingan berurutan dalam jumlah besar terbukti menimbulkan lupa yang bertahap sampai katastrofik~\cite{gupta2024forgetting} dan merusak kemampuan umum model~\cite{gu2024harms}. KoLEG melaporkan kinerja yang tetap baik pada penyuntingan berurutan pengetahuan hukum Korea dengan bantuan temu kembali berkelanjutan~\cite{seo2025koleg}, tetapi pengetahuan yang disunting tetap berada di model, bukan di basis aturan yang dapat diperiksa. Pada konsolidasi undang-undang dengan LLM, kemiripan teks hasilnya mencapai 93--99\%, sedangkan kecocokan persisnya hanya 50,3\% untuk amandemen satu langkah dan 20,51\% untuk rangkaian amandemen~\cite{prior2025risks}. Pada kode, CodeUpdateArena dan ReCode menguji apakah LLM dapat memakai versi API yang sudah berubah~\cite{liu2024codeupdatearena,wu2026recode}. Evaluasi berbasis eksekusi seperti \en{pass@k} dan \en{n@k} menilai program buatan LLM dari hasil menjalankannya, bukan dari kemiripan teksnya~\cite{chen2021evaluating,li2022alphacode}. Prediksi selektif dengan kurva risiko terhadap cakupan memberi kerangka untuk sistem yang boleh menolak menjawab~\cite{geifman2017selective}.

Kelompok ini menjadi dasar metrik RQ2 dan RQ3. SER dan CAR adalah bentuk risiko dan cakupan pada prediksi selektif, konsensus perilaku memakai gagasan penilaian berbasis eksekusi, dan hasil konsolidasi undang-undang di atas menunjukkan mengapa kebenaran tidak boleh diukur dari kemiripan teks. Celahnya, metrik-metrik ini belum diterapkan pada basis aturan hukum yang dieksekusi dan berversi waktu. Penyuntingan pengetahuan juga mengubah bobot model, sedangkan penelitian ini menyunting \en{knowledge base} eksternal yang dapat diperiksa dan dikembalikan.

\subsection{Representasi temporal dan \en{knowledge graph} hukum}
\label{subsec:kg-hukum}

De Martim memodelkan norma hukum dalam \en{knowledge graph} berbasis ontologi yang hierarkis dan temporal, dengan ekspresi berversi agar jawaban tidak memakai versi hukum yang salah~\cite{demartim2025satgraph}. Standar ELI, OWL-Time, PROV-O, SHACL, dan LegalRuleML menyediakan kosakata untuk identitas peraturan, waktu, asal-usul, batasan graf, dan aturan hukum~\cite{publicationsoffice2026eli,cox2017owltime,lebo2013provo,knublauch2017shacl,palmirani2021legalruleml}. \en{Benchmark} 2026 menunjukkan bahwa LLM sering menerapkan ketentuan yang sudah diganti amandemen~\cite{prior2026asking}, salah memilih versi hukum yang berlaku~\cite{huang2026wrong}, dan bahwa sistem RAG statis gagal mengambil versi pasal pajak yang berlaku pada tanggal yang ditanyakan~\cite{cymbler2026temporal}. Rodrigues dkk. memakai pola yang mirip dengan penelitian ini, yaitu LLM mengekstraksi peristiwa perubahan dan gerbang deterministik memutuskan penerimaannya, tetapi untuk memperbarui indeks temu kembali, bukan aturan yang dieksekusi~\cite{rodrigues2026retrieving}. Penelitian ini memakai ulang standar-standar tersebut. Bedanya, graf di penelitian ini memuat aturan yang dieksekusi untuk menghasilkan angka pajak, dan versi waktunya dipakai untuk menjaga hasil masa lampau.

\subsection{Konflik norma, verifikasi \en{knowledge base}, dan pengujian perangkat lunak pajak}
\label{subsec:konflik-uji}

Deteksi konflik norma telah dikaji untuk \en{compliance checking} dengan norma yang saling bertentangan~\cite{robaldo2024compliance}, logika normatif berwaktu~\cite{kharraz2024conflicts}, dan klausul kontrak~\cite{khoja2026consistency}. Anomali pada \en{rule base} yang terus diperbarui dapat dideteksi dari relasi dependensi antaraturan~\cite{boujelben2025anomalies}. Pada pengujian perangkat lunak pajak, spesifikasi eksplisit dan acuan kebenaran umumnya tidak tersedia, sehingga \en{metamorphic testing} dipakai untuk menemukan \en{bug}~\cite{tizpaz2023metamorphic,li2025mrgeneration}. \en{Mutation analysis} dipandang sebagai salah satu kriteria kecukupan uji yang terkuat, dan mutan terbukti berkaitan dengan \en{fault} nyata~\cite{petrovic2021mutation,petrovic2022mutation}. Penelitian ini memakai karya-karya tersebut sebagai dasar rancangan \en{ground truth} dan gerbang G3, bukan sebagai kontribusi.

\begin{table*}[t]
\centering
\caption{Posisi penelitian terhadap karya terkait. Tanda -- berarti aspek itu tidak dinyatakan di abstrak atau bukan fokus karya tersebut}
\label{tab:posisi}
\scriptsize
\begin{tabularx}{\textwidth}{@{}P{2.6cm} P{1.9cm} L P{1.9cm} P{1.6cm} P{1.5cm} P{1.4cm} P{2.3cm}@{}}
\toprule
\textbf{Karya} & \textbf{Domain} & \textbf{Representasi} & \textbf{Penyusun pengetahuan} & \textbf{Pembaruan inkremental} & \textbf{Perambatan kesalahan diukur} & \textbf{Lapisan perusahaan} & \textbf{Evaluasi} \\
\midrule
Catala~\cite{merigoux2021catala} & pajak dan tunjangan (AS, Prancis) & DSL dengan \en{default logic} & manusia & -- & -- & -- & menemukan \en{bug} implementasi resmi \\
Dumancic dkk.~\cite{dumancic2021payroll} & payroll (Belgia) & representasi formal ke \en{constraint} & manusia & -- & -- & -- (optimasi kebijakan) & dipakai lembaga jaminan sosial \\
Makowski dan Tr\k{a}bka~\cite{makowski2025dmn} & akuntansi pajak (Swiss) & DMN & AI generatif dan manusia & -- & -- & -- & studi kasus \\
NTOL~\cite{veeramani2026neurosymbolic} & PPh (India) & \en{knowledge graph} statuta dan aturan simbolik & -- & -- & -- & -- & 20 skenario; akurasi 80\% \\
Yadamsuren dkk.~\cite{yadamsuren2026formalization} & IRC \S121 (AS) & Prolog & LLM, lalu disempurnakan & -- & -- & -- & deteksi inkonsistensi \\
Lorenzo dkk.~\cite{lorenzo2025translating} & hukum pajak & Catala & LLM & -- & -- & -- & \en{benchmark} generasi kode \\
Graus~\cite{graus2026executable} & perizinan (Belanda) & model keputusan yang dieksekusi & LLM & -- & -- & -- & kesetaraan fungsional pada skenario uji \\
Law\&Order~\cite{han2026autoformalization} & formulir pajak (AS) & program simbolik & LLM dengan verifikasi dan perbaikan & -- & -- & -- & 100\% pada 51 SPT \en{holdout} \\
Gogani-Khiabani dkk.~\cite{goganikhiabani2025maintaining} & perangkat lunak pajak (AS) & kode program & LLM (usulan) & ya, sebagai gagasan & -- & -- & kajian tantangan teknis \\
\midrule
\textbf{Penelitian ini} & PPh 21 (Indonesia) & \en{knowledge graph} temporal dua lapis dan aturan yang dieksekusi & LLM dengan gerbang & ya & ya & ya & \en{regulation replay} \\
\bottomrule
\end{tabularx}
\end{table*}

\status{Sel karya 2025--2026 diisi dari abstrak. Metadata karya-karya itu diverifikasi lewat Crossref, arXiv, dan ACL Anthology, tetapi naskah lengkapnya belum dibaca. Kolom pembaruan inkremental dan perambatan kesalahan harus dicek ulang dari naskah lengkap sebelum klaim kebaruan dianggap berdiri (Tahap 1), terutama untuk Gogani-Khiabani dkk., Law\&Order, dan NTOL. Beberapa rujukan adalah \en{preprint} (Law\&Order, Huang dkk., Gogani-Khiabani dkk., CodeUpdateArena, Chen dkk.) atau terbit sebelum 2021. Keduanya bergantung pada keputusan terbuka tentang kriteria rujukan.}

% ============================================================================
\section{Representasi pengetahuan dan \en{inference engine}}
\label{sec:representasi}
% ============================================================================

Bagian ini menetapkan bentuk final \en{knowledge base}, cara konflik antarlapis diselesaikan, dan cara \en{inference engine} menghitung. Bagian ini menjadi dasar kontribusi K1.

\status{Rancangan. \en{Knowledge graph} belum dibangun. Prototipe menyimpan aturan dalam berkas YAML yang divalidasi skema JSON (\texttt{kb/regulasi/}, \texttt{kb/perusahaan/}) dan sudah menjalankan langkah inferensi di Algoritma~\ref{alg:inferensi}, kecuali pembacaan aturan dari graf. Spesifikasi ontologi ditulis di Tahap 3; migrasi YAML ke graf dan pemeriksaan kesetaraan hasil dilakukan di Tahap 5.}

\subsection{Model pengetahuan}
\label{subsec:model}

Pengetahuan PPh 21 terdiri atas dua jenis yang sifatnya berbeda (Tabel~\ref{tab:pembagian}). Pengetahuan struktural menyatakan hubungan. Contohnya regulasi memuat pasal, aturan bersumber dari pasal, PMK baru mengubah atau mencabut PMK lama, uang lembur digolongkan sebagai penghasilan teratur, dan kebijakan perusahaan berkonflik dengan ketentuan wajib. Pengetahuan komputasional menyatakan cara menghitung. Contohnya TER dikali bruto, pembulatan PKP, \en{gross-up} sebagai \en{fixed point}, dan perhitungan ulang di masa pajak terakhir.

\begin{table}[t]
\centering
\caption{Pembagian peran \en{knowledge graph} dan \en{inference engine}}
\label{tab:pembagian}
\scriptsize
\begin{tabularx}{\columnwidth}{@{}P{1.85cm} L P{1.9cm}@{}}
\toprule
\textbf{Jenis} & \textbf{Contoh} & \textbf{Ditangani oleh} \\
\midrule
Struktural & Regulasi memuat pasal; aturan bersumber dari pasal; PMK baru mengubah atau mencabut PMK lama; lembur digolongkan penghasilan teratur; kebijakan berkonflik dengan aturan wajib; aturan membaca dan menghasilkan fakta & \en{Knowledge graph} \\
Komputasional & TER $\times$ bruto; pembulatan per pasal; \en{gross-up} sebagai \en{fixed point}; perhitungan ulang di masa pajak terakhir & \en{Inference engine}; ekspresinya disimpan sebagai atribut \en{node} aturan \\
\bottomrule
\end{tabularx}
\end{table}

Penelitian ini menyimpan pengetahuan struktural sebagai \en{legal knowledge graph} berbasis RDF/OWL. Ekspresi perhitungan disimpan sebagai atribut \en{node} aturan di graf yang sama, lalu \en{inference engine} deterministik membaca aturan dari graf dan menjalankan ekspresinya. Perhitungan tidak dipindahkan ke penalaran ontologi. OWL dan SWRL tidak menyediakan aritmetika eksak dengan aturan pembulatan per pasal, dan tidak menyediakan iterasi \en{fixed point}. Padahal prototipe \en{engine} yang dipakai ulang di penelitian ini sudah tidak menghasilkan selisih yang tidak terjelaskan pada 1.213 nilai harapan dari 45 contoh perhitungan resmi di PER-16/PJ/2016, PP 58/2023, PMK 168/2023, dan ketiga PMK DTP. Dari nilai-nilai itu, 1.185 cocok langsung, 20 berbeda karena kekeliruan contoh resmi yang dikoreksi lewat protokol erratum, dan 8 bergantung pada tafsir atas ketentuan yang tidak tegas. Pembagian peran ini mempertahankan ketepatan tersebut. Seluruh eksperimen prototipe harus menghasilkan keluaran yang identik sebelum dan sesudah migrasi ke graf.

\begin{definisi}[\en{Knowledge base}]
\label{def:kb}
\en{Knowledge base} adalah $\K = \langle \mathcal{O}, \mathcal{R}, \mathcal{P}, \mathcal{T} \rangle$, dengan $\mathcal{O}$ ontologi beserta individunya (regulasi, pasal, dokumen kebijakan, konsep, dan fakta), $\mathcal{R}$ himpunan aturan, $\mathcal{P}$ himpunan parameter berversi, dan $\mathcal{T}$ himpunan tabel berversi. Setiap unsur $\mathcal{R}$, $\mathcal{P}$, dan $\mathcal{T}$ berada pada satu lapisan $\ell \in \{\mathit{reg}, \mathit{prs}\}$, yaitu regulasi atau perusahaan. Lapisan perusahaan disimpan sebagai \en{named graph} tersendiri untuk setiap perusahaan, sehingga beberapa perusahaan dapat memakai lapisan regulasi yang sama.
\end{definisi}

\begin{definisi}[Aturan]
\label{def:aturan}
Aturan adalah tupel
\[ r = \langle \mathit{id}, \ell, \sigma, [t_s, t_e], t_c, \varphi, f \coloneqq e, \mathit{src} \rangle. \]
Sifat $\sigma \in \{\mathit{wajib}, \mathit{tafsir}, \mathit{default}, \mathit{opsional}\}$ menentukan siapa boleh menimpa siapa. Aturan wajib tidak dapat ditimpa lapisan perusahaan. Aturan tafsir menyatakan pilihan peneliti ketika regulasi tidak tegas, beserta rentang tafsir lainnya. Aturan \en{default} boleh ditimpa aturan opsional perusahaan. Interval $[t_s, t_e]$ adalah masa berlaku (\en{valid time}), dengan $t_e = \infty$ bila aturan belum diakhiri, dan $t_c$ adalah tanggal aturan dicatat ke \en{knowledge base} (\en{knowledge time}). Kondisi $\varphi$ dievaluasi atas fakta dan data masukan. Bagian $f \coloneqq e$ menyatakan bahwa aturan menghasilkan fakta $f$ dengan ekspresi $e$. Himpunan $\mathit{src}$ berisi pasal atau butir sumber, dan tidak boleh kosong untuk aturan lapisan regulasi.
\end{definisi}

\begin{definisi}[Hasil perhitungan]
\label{def:hasil}
$\hasil(\K, x, t, \omega)$ adalah nilai fakta keluaran untuk kasus $x$ pada masa pajak $t$, yang dihitung hanya dengan unsur $\K$ yang tanggal pencatatannya $t_c \le \omega$. Fakta keluaran utama adalah PPh 21 yang dipotong per pegawai per masa. Fakta pendukungnya adalah PPh setahun, bruto, PKP, dan \en{take home pay}. Semua fakta keluaran bernilai rupiah bulat.
\end{definisi}

\status{Sifat \emph{tafsir} tidak ada di daftar \en{skeleton} (wajib, \en{default}, opsional), tetapi sudah dipakai prototipe dan catatan revisi. Draf ini mengusulkannya sebagai sifat keempat.}

\subsection{Ontologi dan standar yang dipakai ulang}
\label{subsec:ontologi}

Gambar~\ref{fig:ontologi} menunjukkan kelas dan relasi utama. Kosakatanya memakai ulang standar yang sudah ada, bukan istilah buatan sendiri. ELI (\en{European Legislation Identifier}) dipakai untuk regulasi, bagian-bagiannya, dan relasi amandemen~\cite{publicationsoffice2026eli}. OWL-Time dipakai untuk masa berlaku~\cite{cox2017owltime}, dan PROV-O untuk asal-usul serta waktu pencatatan~\cite{lebo2013provo}. Dengan PROV-O, setiap unsur yang ditambahkan alur akuisisi juga menyimpan dokumen sumber, kutipan pendukung, dan identitas pembaruan yang menghasilkannya.

\begin{figure*}[t]
\centering
\begin{tikzpicture}[
  font=\scriptsize,
  kelas/.style={draw=black!70, rounded corners=2pt, align=center, text width=28mm, minimum height=9mm, fill=white},
  panah/.style={-{Stealth[length=1.8mm]}, semithick},
  ket/.style={font=\tiny, align=center, fill=white, inner sep=1pt}
]
\node[kelas, fill=lapreg!25] (reg) at (0,0) {Regulasi\\(UU, PP, PMK, PER)};
\node[kelas] (pasal) at (0,-1.7) {Pasal atau butir};
\node[kelas, fill=lapprs!30] (keb) at (0,-3.4) {Kebijakan perusahaan\\(PP, PKB, SK direksi)};
\node[kelas, text width=44mm, minimum height=14mm, fill=gray!8] (atr) at (6.2,-1.7) {\textbf{Aturan}\\lapisan $\ell$, sifat $\sigma$, kondisi $\varphi$,\\ekspresi $f \coloneqq e$};
\node[kelas] (kon) at (12.2,0) {Konsep\\(fakta, jenis penghasilan,\\kategori pajak)};
\node[kelas] (par) at (12.2,-1.7) {Parameter berversi};
\node[kelas] (tab) at (12.2,-3.4) {Tabel berversi};
\node[kelas, text width=62mm, fill=violet!6] (ver) at (6.2,-4.5) {Versi: masa berlaku $[t_s, t_e]$ (OWL-Time),\\tanggal pencatatan $t_c$ dan asal-usul (PROV-O)};
\draw[panah] (reg) -- node[ket, right]{memuat} (pasal);
\draw[panah] (keb) -- node[ket, right]{memuat} (pasal);
\draw (reg) edge[panah, loop above, looseness=5] node[ket, above]{mengubah, mencabut} (reg);
\draw[panah] (atr) -- node[ket, above]{bersumber} (pasal);
\draw[panah] (atr) -- node[ket, above, sloped]{menghasilkan,\\membaca} (kon);
\draw[panah] (atr) -- node[ket, above]{membaca} (par);
\draw[panah] (atr) -- node[ket, below, sloped]{membaca} (tab);
\draw (atr) edge[panah, loop above, looseness=5] node[ket, above]{berkonflik dengan, menggantikan} (atr);
\draw[panah] (atr) -- node[ket, right]{berversi} (ver);
\draw[panah] (par.south west) -- (ver.north east);
\draw[panah] (tab) -- (ver);
\end{tikzpicture}
\caption{Ontologi \en{knowledge base}. Kebijakan perusahaan dimodelkan setara dengan regulasi. Yang membedakan keduanya hanya lapisan dan otoritas penerbitnya. Warna biru dan oranye menandai lapisan regulasi dan lapisan perusahaan.}
\label{fig:ontologi}
\end{figure*}

Kebijakan perusahaan dimodelkan setara dengan regulasi pemerintah (Tabel~\ref{tab:setara}). Isi lapisan perusahaan meliputi aturan kebijakan seperti prorata hari kerja, kenaikan gaji di tengah bulan, THR, BPJS, dan \en{take home pay}; pemetaan komponen gaji ke kategori pajak; isian data tambahan yang diminta kebijakan baru; dan aturan pembulatan perusahaan. Dengan pemodelan ini, hubungan antara kebijakan dan regulasi menjadi relasi eksplisit di graf. Akibatnya pertanyaan lintas lapis, misalnya kebijakan perusahaan mana yang berkonflik setelah sebuah regulasi terbit, dapat dijawab dengan satu \en{query}.

\begin{table}[t]
\centering
\caption{Regulasi pemerintah dan kebijakan perusahaan di dalam graf}
\label{tab:setara}
\scriptsize
\begin{tabularx}{\columnwidth}{@{}P{1.5cm} L L@{}}
\toprule
\textbf{Unsur} & \textbf{Regulasi pemerintah} & \textbf{Kebijakan perusahaan} \\
\midrule
Otoritas & Kemenkeu, DJP & Perusahaan \\
Dokumen sumber & UU, PP, PMK, PER & Peraturan perusahaan, PKB, SK direksi \\
Unit sumber & Pasal & Pasal atau butir dokumen internal \\
Sifat aturan & Wajib, tafsir, \en{default} & Opsional, menimpa \en{default} yang sah \\
Versi waktu & Masa berlaku dan tanggal pencatatan & Masa berlaku dan tanggal pencatatan \\
\bottomrule
\end{tabularx}
\end{table}

SHACL menjadi gerbang validasi struktural, termasuk untuk pembaruan yang diusulkan LLM~\cite{knublauch2017shacl}. \en{Shape} dengan tingkat \texttt{sh:Violation} menolak pembaruan, misalnya aturan lapisan regulasi yang tidak bersumber dari pasal, atau aturan pengganti yang mulai berlaku sebelum aturan yang digantikannya. \en{Shape} dengan tingkat \texttt{sh:Warning} tidak menolak, tetapi mencatat konflik antara kebijakan perusahaan dan ketentuan wajib. Konflik seperti ini memang harus tercatat, karena dokumen kebijakan yang sah pun dapat bertentangan dengan regulasi. SPARQL dipakai untuk analisis dampak lintas lapis. Gambar~\ref{lst:rdf} memberi contoh ketiganya untuk konflik klasifikasi lembur di Bagian~\ref{subsec:kebijakan}.

\begin{figure}[t]
\begin{lstlisting}[language=rdf]
# Data: aturan wajib dan aturan Perusahaan X
reg:KLS-LEMBUR a kb:Aturan ; kb:lapisan kb:Regulasi ;
  kb:sifat kb:Wajib ; kb:menggolongkan kb:Lembur ;
  kb:keKategori kb:Teratur ;
  kb:bersumber reg:PER16-Ps1-a15 , reg:PMK168-Ps5-3a .
px:PX-LEMBUR a kb:Aturan ; kb:lapisan kb:Perusahaan ;
  kb:sifat kb:Opsional ; kb:menggolongkan kb:Lembur ;
  kb:keKategori kb:TidakTeratur ; kb:bersumber px:PP-Ps12 .

# SHACL: catat konflik dengan klasifikasi wajib
kb:KlasifikasiWajib a sh:NodeShape ;
  sh:targetClass kb:Aturan ; sh:severity sh:Warning ;
  sh:sparql [ sh:select """SELECT $this WHERE {
    $this kb:lapisan kb:Perusahaan ;
          kb:menggolongkan ?j ; kb:keKategori ?k1 .
    ?w kb:sifat kb:Wajib ; kb:menggolongkan ?j ;
       kb:keKategori ?k2 . FILTER (?k1 != ?k2) }""" ] .

# SPARQL: kebijakan yang menyentuh fakta dari regulasi R
SELECT DISTINCT ?kebijakan WHERE {
  ?a kb:bersumber/kb:bagianDari reg:PMK-168-2023 ;
     (kb:menghasilkan|kb:menggolongkan) ?x .
  ?kebijakan kb:lapisan kb:Perusahaan ;
     (kb:membaca|kb:menghasilkan|kb:menggolongkan) ?x . }
\end{lstlisting}
\caption{Ilustrasi data, \en{shape} SHACL, dan \en{query} SPARQL. Nama kelas dan properti masih usulan dan akan ditetapkan di spesifikasi ontologi.}
\label{lst:rdf}
\end{figure}

\definisikan{(1) Daftar kelas, properti, dan pemetaannya ke ELI, OWL-Time, dan PROV-O. (2) Daftar lengkap \en{shape} SHACL beserta tingkatnya. (3) Bentuk kanonik unsur, yaitu normalisasi aturan agar unsur hasil LLM dapat dicocokkan dengan unsur emas pada metrik diagnostik.}

\subsection{Bahasa ekspresi dan semantik temporal}
\label{subsec:ekspresi}

Ekspresi $e$ dan kondisi $\varphi$ ditulis dalam bahasa ekspresi kecil yang divalidasi sebagai pohon sintaks. Bahasa ini hanya mengizinkan operator aritmetika dan perbandingan, rujukan ke fakta, dan himpunan fungsi bawaan yang terbatas. Prototipe menyediakan 33 fungsi bawaan, misalnya \texttt{tarif\_ter}, \texttt{pasal17}, \texttt{ptkp}, \texttt{bulatkan}, dan fungsi tanggal. Batasan ini juga membatasi keluaran LLM, karena ekspresi yang memanggil fungsi di luar daftar langsung ditolak.

Masa berlaku $[t_s, t_e]$ menentukan aturan mana yang berlaku pada masa pajak $t$, sedangkan $t_c$ menentukan apa yang diketahui sistem pada tanggal $\omega$. Kedua sumbu ini terpisah. Akibatnya regulasi yang berlaku surut, misalnya PMK 10/2025 yang terbit 4 Februari 2025 tetapi berlaku untuk masa Januari 2025, dapat dinyatakan dengan $t_s < t_c$. Pertanyaan \emph{berapa yang seharusnya dipotong untuk Januari menurut hukum yang kini diketahui} dan \emph{berapa yang dihitung sistem pada akhir Januari} keduanya dapat dijawab.

\definisikan{(1) Daftar fungsi yang diizinkan beserta sistem tipenya. (2) Batas ekspresivitas, yaitu jenis ketentuan yang sengaja tidak dapat dinyatakan. Beberapa fungsi bawaan prototipe, seperti \texttt{tarif\_ter} dan \texttt{dtp\_sektor\_tahun}, memuat struktur regulasi di dalam kode \en{engine}. Fungsi seperti ini dapat memaksa perubahan kode ketika regulasi berubah strukturnya, misalnya bila muncul kategori TER keempat, sehingga harus dikaji ulang terhadap syarat A1 sebelum E1.}

\subsection{Resolusi konflik antarlapis}
\label{subsec:resolusi}

Bila lebih dari satu aturan menghasilkan fakta yang sama pada masa yang sama, \en{engine} memilih satu pemenang. Misalkan $C(f, t, \omega)$ adalah himpunan aturan yang menghasilkan fakta $f$, berlaku pada masa $t$, sudah dicatat pada tanggal $\omega$, dan kondisinya terpenuhi. Aturan yang dicabut secara eksplisit lewat relasi \emph{menggantikan} sudah dikeluarkan dari $C$. \en{Engine} lalu membandingkan kunci leksikografis
\begin{equation}
\kappa(r) = \big(\mathit{sup}(r),\ \mathit{ovr}(r),\ \mathit{spec}(r),\ t_s(r),\ \mathit{prio}(r)\big).
\label{eq:kunci}
\end{equation}
Komponen $\mathit{sup}(r)$ bernilai 1 bila $r$ aturan wajib atau tafsir di lapisan regulasi, sesuai asas \en{lex superior}. Komponen $\mathit{ovr}(r)$ bernilai 1 bila $r$ aturan opsional perusahaan, sehingga aturan perusahaan mengalahkan aturan \en{default} regulasi. Komponen $\mathit{spec}(r)$ mengukur kekhususan kondisi sesuai asas \en{lex specialis}, $t_s(r)$ mewakili asas \en{lex posterior}, dan $\mathit{prio}(r)$ adalah prioritas eksplisit. Bila lebih dari satu aturan memiliki kunci terbesar dengan nilai berbeda, \en{engine} berhenti dengan status AMBIGU dan tidak memilih diam-diam. Setiap kali aturan wajib mengalahkan aturan perusahaan, \en{engine} mencatat konflik itu di jejak. Untuk contoh lembur, kategori dari regulasi yang dipakai, dan konflik dengan aturan Perusahaan X tercatat beserta pasal sumbernya.

\terbuka{Ukuran \en{lex specialis}. Prototipe menghitung $\mathit{spec}(r)$ dari jumlah konjungsi di kondisi $\varphi$. Alternatifnya adalah relasi superioritas eksplisit antaraturan yang dicatat di graf. Pilihan ini juga menentukan apa yang harus dihasilkan LLM.}

\subsection{\en{Inference engine}}
\label{subsec:engine}

Algoritma~\ref{alg:inferensi} merangkum langkah inferensi untuk satu pegawai dan satu tahun pajak. \en{Engine} memilih aturan menurut kedua sumbu waktu, membangun graf dependensi antarfakta, lalu mengurutkan komponen terhubung kuatnya secara topologis. Komponen yang bersiklus hanya diizinkan bila memuat aturan yang ditandai sebagai \en{fixed point}, seperti \en{gross-up}. Untuk \en{gross-up}, tunjangan pajak $T$ memenuhi $T = g(T)$ dengan $g(T) = \mathrm{PPh}(B + T)$ untuk bruto $B$ tanpa tunjangan. Fungsi $g$ tidak turun, karena PPh tidak pernah turun ketika bruto naik. Pada domain rupiah bulat $\{0, \dots, B\}$, $g$ juga memetakan domain ke dirinya sendiri, karena PPh tidak pernah melebihi setengah bruto. Menurut teorema Knaster--Tarski~\cite{tarski1955lattice}, titik tetapnya ada dan memiliki unsur terkecil dan terbesar. Iterasi dari batas bawah dan batas atas mencapai keduanya dalam jumlah langkah terbatas. Bila keduanya berbeda, \en{engine} mencatat solusi ganda dan memilih secara deterministik, dengan bawaan solusi terkecil.

\begin{algorithm}[t]
\caption{Inferensi untuk satu pegawai dan satu tahun pajak}
\label{alg:inferensi}
\footnotesize
\begin{algorithmic}[1]
\Require $\K$, data pegawai $x$, tahun pajak $y$, tanggal pengetahuan $\omega$
\Ensure nilai setiap fakta per masa beserta jejaknya
\State $R \gets$ aturan dengan $t_c \le \omega$ yang masa berlakunya beririsan dengan $y$
\State keluarkan dari $R$ aturan yang digantikan aturan lain di $R$
\State bangun graf dependensi $G$: $f' \to f$ bila aturan untuk $f$ membaca $f'$
\State hitung komponen terhubung kuat $G$, lalu urutkan topologis
\ForAll{komponen $Q$ menurut urutan topologis}
  \If{$Q$ bersiklus}
    \State galat bila $Q$ tidak memuat aturan \en{fixed point}
    \State iterasi dari batas bawah dan atas; catat bila hasilnya ganda
  \Else
    \ForAll{instans $(f, t)$ dengan $f \in Q$ dan $t$ masa dalam $y$}
      \State pilih pemenang dari $C(f, t, \omega)$ dengan kunci $\kappa$
      \State evaluasi ekspresinya dengan aritmetika eksak
      \State catat nilai, pemenang, pasal, dan aturan yang ditolak
    \EndFor
  \EndIf
\EndFor
\State \Return nilai fakta dan jejak
\end{algorithmic}
\end{algorithm}

Ketepatan sampai rupiah dijaga tiga mekanisme. Pertama, nilai rupiah disimpan sebagai bilangan bulat dan tarif sebagai bilangan rasional, sehingga tidak ada bilangan titik mengambang. Pemindai pohon sintaks menolak kode \en{engine} yang memanggil konversi ke bilangan titik mengambang atau pembulatan bawaan bahasa pemrograman. Kedua, setiap titik pembulatan merujuk ke satu entri registri pembulatan yang mencatat satuan, salah satu dari enam mode pembulatan, dan dasar hukumnya. Ketiga, masa pajak suatu penghasilan ditentukan dari tiga tanggal, yaitu periode kerja, tanggal terutang, dan tanggal dibayar, karena PPh 21 terutang pada saat penghasilan dibayar atau terutang, mana yang lebih dahulu. Setiap nilai keluaran disertai jejak yang memuat aturan pemenang, pasal sumbernya, dan aturan yang ditolak beserta alasannya.

\status{Algoritma~\ref{alg:inferensi}, aritmetika eksak, registri pembulatan, model tiga tanggal, dan jejak sudah ada di prototipe (\texttt{engine/inferensi.py}, \texttt{engine/pembulatan.py}). Yang baru adalah pembacaan aturan dari graf dan uji unit \en{engine} yang terpisah dari \en{knowledge base} (syarat E0).}

% ============================================================================
\section{Akuisisi pengetahuan inkremental berbantuan LLM}
\label{sec:akuisisi}
% ============================================================================

Bagian ini mendefinisikan pembaruan dan adaptivitas secara operasional, lalu menjelaskan alur akuisisi yang menjaganya. Bagian ini menjadi dasar kontribusi K2 dan K3.

\status{Rancangan. Prototipe \texttt{asisten\_kb/} baru dapat mengubah satu PDF peraturan menjadi rancangan berkas aturan lewat satu panggilan LLM dengan keluaran terstruktur, lalu memvalidasi bentuknya dengan \en{engine}. Prototipe ini baru diuji satu kali pada PDF peraturan fiktif. Penilaian relevansi, $m$ kandidat, gerbang G2 dan G5, generator $\SG$, konsensus perilaku, \en{abstain}, dan pencatat eksperimen belum dibangun (Tahap 9).}

\subsection{Pembaruan dan adaptivitas}
\label{subsec:adaptif}

\begin{definisi}[Peristiwa perubahan]
\label{def:peristiwa}
Peristiwa perubahan adalah $\epsilon = \langle d, t_{\mathit{terbit}}, t_{\mathit{berlaku}}, \tau \rangle$, yang terdiri atas dokumen sumber $d$, tanggal terbit, tanggal mulai berlaku, dan jenis perubahan $\tau$ menurut Tabel~\ref{tab:jenis}. Peristiwa berlaku surut bila $t_{\mathit{berlaku}} < t_{\mathit{terbit}}$. Rangkaian $\epsilon_1, \dots, \epsilon_N$ diurutkan menurut $t_{\mathit{terbit}}$.
\end{definisi}

\begin{definisi}[Pembaruan]
\label{def:pembaruan}
Pembaruan adalah $\Delta = \langle \Delta^{+}, \Delta^{\sim}, \Delta^{-} \rangle$, yang berisi unsur yang ditambah, diubah, dan dicabut, dengan $\K' = \K \oplus \Delta$. Pembaruan emas $\Delta^{\emas}_n$ disusun peneliti untuk peristiwa $\epsilon_n$ dan memenuhi $\K^{\emas}_n = \K^{\emas}_{n-1} \oplus \Delta^{\emas}_n$. Pembaruan dari alur akuisisi ditulis $\hat{\Delta}_n$. Pembaruan kosong ($\Delta = \langle \emptyset, \emptyset, \emptyset \rangle$) adalah jawaban yang benar untuk dokumen yang tidak mengubah perhitungan.
\end{definisi}

\begin{definisi}[Syarat adaptif]
\label{def:syarat}
Untuk $\K' = \K \oplus \Delta$ yang menanggapi peristiwa $\epsilon$, dan himpunan kasus $S$, syaratnya adalah sebagai berikut.
\begin{itemize}[leftmargin=3.4em,labelsep=0.4em,nosep]
  \item[(A1)] Invariansi kode. Kode \en{engine} tidak berubah.
  \item[(A2)] Kesesuaian. $\hasil(\K', x, t, \omega) = \hasil(\K^{\emas}, x, t, \omega)$ untuk setiap $x \in S$, $t \geq t_{\mathit{berlaku}}$, dan $\omega \geq t_{\mathit{terbit}}$.
  \item[(A3$_t$)] Isolasi masa. Hasil untuk $t < t_{\mathit{berlaku}}$ tidak berubah, untuk semua $\omega$.
  \item[(A3$_\omega$)] Isolasi pengetahuan. Hasil untuk $\omega < t_{\mathit{terbit}}$ tidak berubah.
  \item[(A4)] Konsistensi lintas lapis. Setiap kebijakan perusahaan yang berkonflik dengan aturan wajib di $\K'$ dilaporkan.
\end{itemize}
\end{definisi}

\begin{definisi}[Adaptivitas representasional]
\label{def:adapt-rep}
Representasi adaptif terhadap $\epsilon$ bila \emph{ada} $\Delta$ dalam bahasa representasi yang memenuhi A1--A4. Ini sifat representasi. Peneliti membuktikannya untuk setiap peristiwa dengan menyusun $\Delta^{\emas}$ (RQ1).
\end{definisi}

\begin{definisi}[Adaptivitas operasional]
\label{def:adapt-op}
Sistem adaptif secara operasional terhadap $\epsilon$ bila pembaruan $\hat{\Delta}$ yang \emph{dipilih alur akuisisi} memenuhi A2, A3$_t$, dan A4 pada himpunan kasus evaluasi $\SE$. Ini sifat alur akuisisi dan diukur sebagai laju atas peristiwa (RQ2, RQ3).
\end{definisi}

Dua syarat berlaku karena arsitekturnya. A1 berlaku karena LLM hanya menyunting \en{knowledge base} dan tidak pernah menyentuh kode \en{engine}. A3$_\omega$ berlaku karena konstruksi bila setiap unsur baru dicatat dengan $t_c = t_{\mathit{terbit}}$, sebab $\hasil$ untuk $\omega < t_c$ tidak melihat unsur itu. Karena itu kedua syarat diperiksa sebagai invarian, bukan diuji sebagai hipotesis. Syarat yang bergantung pada LLM, yaitu A2, A3$_t$, dan A4, adalah yang diuji.

\definisikan{Semantik berlaku surut menurut mekanisme hukumnya. Untuk PMK 10/2025, harus ditetapkan bagaimana PPh 21 masa Januari 2025, yang dipotong sebelum PMK itu terbit, diselesaikan, misalnya lewat pembayaran tunai DTP atau perhitungan ulang di masa pajak terakhir, bukan dengan menghitung ulang potongan bulan yang sudah lewat. Nilai \emph{sebagaimana seharusnya} untuk $t < t_{\mathit{terbit}}$ mengikuti mekanisme itu. Tanpa definisi ini, \en{ground truth} untuk peristiwa berlaku surut tidak terdefinisi.}

\subsection{Alur akuisisi}
\label{subsec:alur}

Gambar~\ref{fig:alur} menunjukkan delapan tahap alur akuisisi. Pertama, sistem memotong dokumen per pasal. Kedua, LLM menilai apakah dokumen mengubah perhitungan PPh 21 pegawai tetap, dan menyertakan kutipan teks pendukung. Bila dokumen dinilai tidak relevan dan kutipannya lolos G2, sistem menerapkan pembaruan kosong dan berhenti. Ketiga, sistem menyusun konteks dari bagian graf yang relevan, yaitu aturan, fakta, dan parameter yang disebut atau disentuh dokumen. Keempat, LLM menyusun $m$ kandidat $\Delta$ secara independen, dan setiap unsurnya disertai kutipan sumber. Kelima, setiap kandidat melewati gerbang G1--G5. Kandidat yang gagal mendapat umpan balik dari gerbang dan boleh diperbaiki dalam jumlah putaran terbatas. Keenam, sistem memilih satu kandidat lewat konsensus perilaku, atau \en{abstain} bila tidak ada kandidat yang lolos. Ketujuh, peninjau manusia memeriksa pembaruan dengan bantuan laporan dampak. Kedelapan, pembaruan berlaku sebagai versi baru \en{knowledge base}.

\begin{figure}[t]
\centering
\begin{tikzpicture}[
  font=\scriptsize,
  node distance=3mm,
  tahap/.style={draw=black!70, rounded corners=2pt, align=left, text width=44mm, minimum height=6mm, fill=white},
  llm/.style={tahap, fill=violet!7},
  grb/.style={tahap, fill=teal!8},
  sisi/.style={draw=black!50, dashed, rounded corners=2pt, align=center, text width=15mm, minimum height=6mm, fill=gray!6},
  panah/.style={-{Stealth[length=1.6mm]}, semithick},
  ket/.style={font=\tiny, align=center, fill=white, inner sep=1pt}
]
\node[tahap] (s1) {(1) Potong dokumen per pasal};
\node[llm, below=of s1] (s2) {(2) Penilaian relevansi berkutipan};
\node[tahap, below=of s2] (s3) {(3) Susun konteks dari graf};
\node[llm, below=of s3] (s4) {(4) Susun $m$ kandidat $\Delta$ berkutipan};
\node[grb, below=of s4] (s5) {(5) Gerbang G1--G5};
\node[grb, below=of s5] (s6) {(6) Konsensus perilaku pada $\SG$};
\node[tahap, below=of s6] (s7) {(7) Peninjauan dengan laporan dampak};
\node[tahap, below=of s7, fill=green!6] (s8) {(8) Terapkan sebagai versi baru $\K'$};
\node[sisi, right=10mm of s2] (kosong) {$\Delta$ kosong};
\node[sisi, right=10mm of s6] (abs) {\en{Abstain}};
\draw[panah] (s1) -- (s2);
\draw[panah] (s2) -- (s3);
\draw[panah] (s3) -- (s4);
\draw[panah] (s4) -- (s5);
\draw[panah] (s5) -- (s6);
\draw[panah] (s6) -- (s7);
\draw[panah] (s7) -- (s8);
\draw[panah] (s2) -- node[ket, above]{tidak\\relevan} (kosong);
\draw[panah] (s6) -- node[ket, above]{tak ada\\yang lolos} (abs);
\draw[panah] (s5.west) -- ++(-3mm,0) |- (s4.west);
\draw[panah] (kosong.east) -- ++(3mm,0) |- (s8.east);
\draw[panah] (abs) |- (s7.east);
\end{tikzpicture}
\caption{Alur akuisisi pengetahuan inkremental. Kandidat yang gagal di gerbang mendapat umpan balik untuk diperbaiki. Dokumen yang membuat sistem \en{abstain} diteruskan ke peninjau tanpa pembaruan otomatis.}
\label{fig:alur}
\end{figure}

\terbuka{Tabel tarif di lampiran regulasi dapat berisi ratusan baris, misalnya tabel TER. Apakah LLM menyalinnya sendiri, atau boleh memanggil alat ekstraksi tabel PDF yang deterministik? Usulan \en{skeleton}: alat diizinkan dan dicatat di pra-registrasi, sehingga kemampuan memahami aturan terpisah dari kemampuan menyalin angka.}

\subsection{Gerbang verifikasi}
\label{subsec:gerbang}

Tabel~\ref{tab:gerbang} memuat lima gerbang beserta syarat yang dijaganya. G1 memeriksa struktur dengan \en{shape} SHACL (Bagian~\ref{subsec:ontologi}). G2 memeriksa bahwa setiap kutipan benar-benar ada di dokumen sumber, sehingga dasar hukum karangan tertolak. G3 memeriksa konsistensi statis atas tipe, dependensi, konflik, dan anomali \en{knowledge base}. G4 adalah uji regresi pada $\SG$. Hasil untuk masa sebelum $t_{\mathit{berlaku}}$ tidak boleh berubah. Peran G4 setara dengan \en{cornerstone case} pada \en{Ripple-Down Rules}. G5 adalah gerbang non-vakuitas.

\begin{table}[t]
\centering
\caption{Gerbang verifikasi dan syarat yang dijaganya}
\label{tab:gerbang}
\scriptsize
\begin{tabularx}{\columnwidth}{@{}P{0.5cm} L P{1.35cm}@{}}
\toprule
\textbf{ID} & \textbf{Pemeriksaan} & \textbf{Syarat} \\
\midrule
G1 & Struktur dan batasan graf (\en{shape} SHACL) & validitas \\
G2 & Kutipan sumber ada di dokumen & keberadaan kutipan \\
G3 & Konsistensi statis atas tipe, dependensi, konflik, dan anomali & A4 \\
G4 & Regresi pada $\SG$; hasil sebelum $t_{\mathit{berlaku}}$ tidak berubah & A3$_t$ \\
G5 & Non-vakuitas untuk $\Delta$ tidak kosong. Setiap unsur yang ditambah atau diubah mengubah minimal satu fakta keluaran pada kasus $\SG$ sesudah $t_{\mathit{berlaku}}$, termasuk fakta selain PPh 21. Setiap unsur yang dicabut tidak lagi menyala. Tanggal dan subjek yang disebut dokumen tercermin di $\Delta$ & menolak pembaruan tanpa efek \\
\bottomrule
\end{tabularx}
\end{table}

G5 diperlukan karena celah pada empat gerbang lainnya. Pembaruan kosong, atau aturan yang tidak pernah menyala, lolos G1--G4, karena keduanya tidak menimbulkan regresi. Putaran perbaikan yang hanya menerima umpan balik regresi bahkan mendorong kandidat ke arah itu. G5 menutup celah ini. G5 tidak bertentangan dengan dokumen pengecoh, karena pembaruan kosong hanya dapat diterapkan lewat penilaian relevansi yang kutipannya lolos G2, bukan lewat gerbang.

Kasus $\SG$ dibangkitkan otomatis dari $\K_{n-1} \oplus \hat{\Delta}$. Generatornya menyasar kondisi $\varphi$ setiap unsur yang berubah, sehingga pengetahuan manusia tentang isi dokumen tidak bocor ke alur. $\SG$ hanya dipakai gerbang, dan tidak beririsan dengan $\SE$ yang dipakai untuk menilai sistem.

\subsection{Pemilihan kandidat dan \en{abstain}}
\label{subsec:konsensus}

Untuk setiap kandidat $\hat{\Delta}_j$ yang lolos semua gerbang, sistem menghitung vektor hasil
\begin{equation}
v_j = \big(\hasil(\K \oplus \hat{\Delta}_j, x, t, \omega)\big)_{(x,t,\omega) \in \SG}.
\label{eq:vektor}
\end{equation}
Kandidat dengan vektor yang sama dikelompokkan, dan sistem memilih wakil kelompok terbesar. Cara ini disebut konsensus perilaku, karena kandidat dikelompokkan menurut hasil perhitungannya, bukan menurut kemiripan bentuk aturannya. Dua kandidat yang ditulis berbeda tetapi selalu memberi angka yang sama dianggap setara. Bila tidak ada kandidat yang lolos, sistem \en{abstain} dan meneruskan dokumen ke peninjau manusia.

\definisikan{(1) Daftar \en{constraint} G1 dan tingkatnya. (2) Cara G2 mencocokkan kutipan, misalnya kecocokan setelah normalisasi spasi dan nomor pasal. (3) Generator $\SG$ dan jaminan bahwa $\SG \cap \SE = \emptyset$. (4) Nilai $m$ (usulan 5) dan batas putaran perbaikan. (5) Aturan bila dua kelompok sama besar. Usulan draf ini: \en{abstain}. (6) Syarat pembekuan. G5 harus menerima $\Delta^{\emas}$ dari seluruh peristiwa pengembangan, artinya laju penolakan kelirunya nol.}

\subsection{Analisis dampak lintas lapis}
\label{subsec:dampak}

Diberikan pembaruan $\Delta$, detektor menelusuri graf secara statis lewat \en{query} SPARQL untuk menemukan kebijakan perusahaan yang terdampak. Sebuah kebijakan disebut terdampak oleh $\Delta$ bila memenuhi minimal satu dari tiga jenis berikut. Jenis (a), nilai fakta keluaran berubah untuk minimal satu kasus. Jenis (b), aturan pemenang untuk fakta yang dibaca atau dihasilkan kebijakan itu berubah. Jenis (c), kebijakan itu menjadi berkonflik dengan aturan wajib. Label jenis (a) dan (b) diturunkan dari eksekusi $\K^{\emas}$ pada $\SE$, sedangkan detektor bekerja statis pada graf tanpa menjalankan \en{engine}. Karena sumber label dan cara kerja detektor berbeda, evaluasinya tidak tautologis. Label jenis (c) diberikan dua pelabel secara independen.

\subsection{Konfigurasi LLM}
\label{subsec:konfigurasi}

Untuk setiap jalan eksperimen, sistem mencatat model dan versinya, tanggal akses, \en{training cutoff} model, parameter pembangkitan, format keluaran terstruktur, isi \en{prompt}, serta seluruh masukan dan keluaran. Dengan catatan ini, semua analisis dapat diulang tanpa memanggil LLM lagi.

\terbuka{Model utama, model kedua (\en{open-weight}), dan anggaran API. Prototipe asisten memakai API OpenAI sebagai bawaan dan juga mendukung API Anthropic. Model kedua hanya dipakai untuk kondisi usulan dan pembanding B3 (Bagian~\ref{subsec:pembanding}).}

% ============================================================================
\section{Metodologi evaluasi}
\label{sec:evaluasi}
% ============================================================================

Evaluasi memakai lima prinsip. Pertama, setiap RQ punya satu metrik utama yang ditetapkan sebelum eksperimen. Kedua, kebenaran diukur dari hasil perhitungan, bukan dari kemiripan bentuk aturan buatan LLM dengan aturan emas. Rata-rata selisih seperti MAE tidak dipakai, karena rata-rata yang kecil dapat menyembunyikan kasus yang salah. Ketiga, data yang dipakai sistem selalu terpisah dari data yang dipakai untuk menilainya. Keempat, alat ukur divalidasi sebelum dipakai mengukur sistem. Kelima, jumlah peristiwa nyata kecil, sehingga hasil dilaporkan sebagai estimasi dengan selang kepercayaan eksak, bukan sebagai keputusan signifikan atau tidak.

\status{Seluruh bagian ini adalah rancangan. Belum ada eksperimen yang dijalankan, dan belum ada pra-registrasi. Ambang hipotesis $\theta_1$--$\theta_4$ dan margin $\delta_S$ masih kosong dan baru dikunci di Tahap 12.}

\subsection{Desain penelitian dan pra-registrasi}
\label{subsec:desain}

Penelitian ini mengikuti \en{design science research}~\cite{hevner2004design,peffers2007design}. Artefaknya adalah KBS adaptif beserta alur akuisisinya, dan evaluasinya adalah rangkaian eksperimen E0--E5. Sebelum eksperimen pada data uji, peneliti mendaftarkan hipotesis, metrik utama, ambang, pembagian data, kasus uji evaluasi, dan rencana analisis ke layanan pra-registrasi pihak ketiga yang memberi stempel waktu, misalnya OSF atau AsPredicted, lalu menautkannya dengan \en{tag} git. Analisis terdiri atas tiga jenis. Pemeriksaan invarian mencakup A1 dan A3$_\omega$, dan pelanggarannya harus nol. Analisis konfirmatori memakai metrik utama untuk menguji hipotesis. Analisis eksploratori memakai metrik diagnostik dan diberi label eksploratori saat dilaporkan.

\subsection{Pembagian data dan kronologi}
\label{subsec:pembagian}

Data dibagi menjadi tiga bagian menurut kronologi.
\begin{itemize}[leftmargin=*,nosep]
  \item \textbf{Pengembangan.} Regulasi fiktif, perubahan historis sebelum rentang penelitian (misalnya perubahan PTKP 2016 dan lapisan Pasal 17 oleh UU HPP), dan sebagian kecil dokumen kebijakan buatan. Peneliti memakainya untuk merancang \en{prompt}, gerbang, dan aturan pemilihan, lalu membekukan ketiganya.
  \item \textbf{Uji.} Seluruh peristiwa nyata pada 2023--2026, termasuk pergantian ke TER. Peneliti tidak melihat peristiwa uji selama mengembangkan alur LLM. Representasi, \en{engine}, dan $\K^{\emas}$ memang disusun dengan mengetahui peristiwa ini, dan akibatnya ditangani di Bagian~\ref{subsec:hindsight}.
  \item \textbf{Uji prospektif.} Setiap regulasi PPh 21 yang terbit sesudah tanggal pra-registrasi. Peristiwa ini bebas kontaminasi data latih LLM dan belum pernah dilihat perancang.
\end{itemize}

\subsection{Bias pandangan belakang}
\label{subsec:hindsight}

Representasi dan \en{engine} prototipe dikembangkan ketika perubahan 2023--2026 sudah diketahui. Karena itu, klaim \emph{tanpa mengubah kode} untuk peristiwa tersebut lemah bila berdiri sendiri. Ada tiga mitigasi. Pertama, peneliti mempublikasikan riwayat pengembangan \en{engine} dan menghitung setiap perubahan \en{engine} yang dipicu suatu regulasi sebagai pelanggaran A1 pada representasi. Kedua, orang lain menulis perubahan struktural sintetis dari taksonomi perubahan, yang disebut perubahan tak terlihat, untuk diuji di RQ1. Ketiga, penelitian ini memakai peristiwa uji prospektif.

Untuk perubahan tak terlihat, penulisnya menyerahkan teks perubahan beserta 5--10 contoh perhitungan yang sudah dihitung, sebelum peneliti melihatnya. Perubahan itu dihitung berhasil, yaitu RA bernilai 1, hanya bila $\Delta$ yang disusun peneliti mereproduksi seluruh contoh tanpa mengubah \en{engine} dalam batas waktu yang ditetapkan. Dengan cara ini peneliti tidak menilai pekerjaannya sendiri. Jumlah perubahan tak terlihat menentukan kekuatan klaim RQ1. Bila delapan perubahan semuanya berhasil, batas bawah eksak RA pada keyakinan 95\% hanya $0{,}05^{1/8} \approx 0{,}69$. Dengan lima belas perubahan, batasnya menjadi sekitar $0{,}82$.

\terbuka{Siapa penulis perubahan tak terlihat, dan berapa jumlahnya? Usulan \en{skeleton}: minimal lima belas.}

\subsection{Rantai kepercayaan dan validasi instrumen}
\label{subsec:instrumen}

Semua angka RQ2--RQ4 bergantung pada $\K^{\emas}$ dan $\SE$. Karena itu keduanya divalidasi lebih dulu sebagai instrumen, bukan sebagai hasil penelitian. Gambar~\ref{fig:rantai} menunjukkan rantainya. Acuan independen memvalidasi $\K^{\emas}_n$, lalu $\K^{\emas}_n$ menghasilkan nilai harapan $\SE$, dan $\SE$ menilai pembaruan dari LLM.

\begin{figure}[t]
\centering
\begin{tikzpicture}[
  font=\scriptsize,
  acuan/.style={draw=black!60, rounded corners=1.5pt, align=center, text width=19mm, minimum height=5mm, fill=gray!8},
  kotak/.style={draw=black!70, rounded corners=2pt, align=center, text width=17mm, minimum height=9mm, fill=white},
  panah/.style={-{Stealth[length=1.6mm]}, semithick},
  ket/.style={font=\tiny, align=center, fill=white, inner sep=1pt}
]
\node[acuan] (g1) at (0,0) {GT1 contoh resmi};
\node[acuan] (g2) at (0,-0.7) {GT2 kalkulator DJP};
\node[acuan] (g3) at (0,-1.4) {GT3 hitung manual};
\node[acuan] (g4) at (0,-2.1) {GT4 pihak ketiga};
\node[acuan, fill=lapprs!25] (g5) at (0,-2.8) {GT5 contoh kebijakan};
\node[kotak, fill=blue!6] (kemas) at (3.1,-1.4) {$\K^{\emas}_n$\\tervalidasi};
\node[kotak, fill=teal!8] (se) at (5.9,-1.4) {$\SE$ dan nilai harapan (kecukupan: \en{mutation score})};
\node[kotak] (nilai) at (5.9,-3.3) {menilai $\hat{\Delta}$\\(A2, A3$_t$, A4)};
\node[kotak, dashed, fill=gray!5] (sg) at (5.9,0.45) {$\SG$, hanya untuk gerbang};
\foreach \g in {g1,g2,g3,g4,g5} \draw[panah] (\g.east) -- (kemas.west);
\draw[panah] (kemas) -- (se);
\draw[panah] (se) -- (nilai);
\draw[dashed] (sg) -- node[ket, right]{tidak beririsan} (se);
\end{tikzpicture}
\caption{Rantai kepercayaan \en{ground truth}. GT1 yang dipakai di sini adalah subset \en{holdout} yang tidak pernah dipakai selama pengembangan.}
\label{fig:rantai}
\end{figure}

\begin{enumerate}[leftmargin=*,nosep]
  \item \textbf{Acuan.} Ada lima sumber kebenaran. GT1 adalah contoh perhitungan resmi di PMK, PP, dan PER, yang dibagi menjadi contoh yang sudah dipakai selama pengembangan dan subset \en{holdout} yang tidak pernah dipakai. GT2 adalah kalkulator otoritas pajak, bila tersedia untuk rezim yang diuji. GT3 adalah perhitungan manual oleh penghitung independen pada sampel bertingkat per lapisan tarif. GT4 adalah implementasi pihak ketiga, yaitu kalkulator \en{open source} dan kalkulator publik, dengan versi dan tanggal yang dicatat. GT5 adalah acuan untuk angka yang ditentukan kebijakan perusahaan, seperti gaji prorata, THR, dan \en{take home pay}. Isinya contoh perhitungan yang ditulis kolaborator dari teks kebijakan sebelum peneliti menyusun $\Delta^{\emas}$, ditambah contoh yang tercantum di dokumen kebijakan itu sendiri.
  \item \textbf{Pelaporan selisih.} Peneliti melaporkan semua selisih menurut hasil adjudikasi, yaitu \en{knowledge base} salah, acuan salah, atau hukum ambigu. Pengadjudikasi kedua memeriksa ulang sampel acak secara buta.
  \item \textbf{Matriks acuan.} Matriks ini menunjukkan sel ruang uji yang punya acuan independen dan sel yang tidak.
  \item \textbf{Uji \en{engine}.} \en{Engine} diuji terpisah dari \en{knowledge base} dengan uji unit terhadap aritmetika yang dihitung tangan, karena $\SE$ juga dihitung dengan \en{engine} yang sama.
  \item \textbf{Kecukupan $\SE$.} Peneliti mengukur \en{mutation score} global dan \en{mutation score} terfokus-$\Delta$ untuk setiap peristiwa, yang hanya memutasi unsur $\Delta^{\emas}_n$~\cite{petrovic2021mutation,petrovic2022mutation}. Kedua skor harus mencapai ambang minimum sebelum E2 dan E3. Kesalahan nyata LLM pada data pengembangan ditambahkan sebagai mutan realistis.
  \item \textbf{Keandalan $\Delta^{\emas}$.} Penyusun pertama menyusun $\Delta^{\emas}$, lalu pembaca kedua yang independen memeriksanya. Kesepakatan atas keputusan unsur dilaporkan dengan Krippendorff's $\alpha$. $\Delta^{\emas}$ dikunci dengan stempel waktu sebelum LLM dijalankan.
  \item \textbf{Syarat lulus E0.} Tidak ada selisih berstatus \en{knowledge base} salah yang belum diselesaikan, dan kedua \en{mutation score} mencapai ambang. E2--E4 baru dijalankan setelah E0 lulus.
  \item \textbf{Protokol koreksi emas.} Bila keluaran LLM menyingkap kesalahan pada $\K^{\emas}$ atau $\Delta^{\emas}$, peneliti memperbaikinya, menilai ulang seluruh kondisi dengan emas yang sudah diperbaiki, dan melaporkan jumlah koreksi. Tanpa protokol ini, emas dapat dicurigai disunting secara selektif.
\end{enumerate}

\status{Bahan awal sudah ada. Prototipe memiliki 45 contoh resmi kanonik dengan 1.213 nilai harapan dan protokol erratum, tetapi seluruhnya sudah dipakai selama pengembangan sehingga tidak dapat menjadi \en{holdout} GT1. Prototipe juga memiliki 52 mutan pada \en{knowledge base} dengan \en{mutation score} 100\%. GT2 (kalkulasi ulang dengan kalkulator DJP) belum pernah dijalankan.}

\terbuka{(1) Apakah ada praktisi pajak bersertifikat, misalnya dari \en{tax center} universitas, yang dapat dibayar beberapa jam untuk memeriksa $\Delta^{\emas}$ peristiwa TER dan sekitar 50 kasus bertingkat? Ini mengurangi kelemahan terbesar penelitian ini, yaitu tidak adanya \en{domain expert}. (2) Siapa pembaca kedua $\Delta^{\emas}$?}

\subsection{Ruang uji}
\label{subsec:ruanguji}

Ruang uji mencakup seluruh lapisan tarif dan kombinasi dimensi di Tabel~\ref{tab:ruang-pajak} dengan tiga kriteria.
\begin{itemize}[leftmargin=*,nosep]
  \item \textbf{Cakupan kombinatorial berbatasan.} Setiap pasangan nilai dari dua dimensi muncul di minimal satu kasus (kombinasi 2-arah), dan cakupan 3-arah dilaporkan sebagai pelengkap. Keduanya dihitung hanya atas kombinasi yang sah menurut batasan antardimensi, misalnya DTP hanya pada 2025--2026.
  \item \textbf{Cakupan batas.} Setiap batas lapisan diuji di tiga titik, yaitu tepat di batas, satu satuan di bawahnya, dan satu satuan di atasnya. Di titik-titik inilah kesalahan pembulatan dan pemilihan lapisan paling sering terjadi. Untuk TER, batasnya adalah bruto bulanan dalam rupiah. Untuk Pasal 17, batasnya ada pada PKP sesudah pembulatan, sehingga generator membangkitkan kasus secara terbalik, dari PKP ke penghasilan bruto.
  \item \textbf{Sampel operasional.} Generator mengambil kasus pegawai-tahun secara acak sederhana dari distribusi upah dan status PTKP menurut data Badan Pusat Statistik. Sampel ini dipakai untuk memperkirakan laju kegagalan dalam pemakaian nyata.
\end{itemize}

Tabel~\ref{tab:cakupan} merangkum jumlah sel yang sudah dapat dihitung. Lima lapisan Pasal 17 memberi empat batas, sehingga untuk empat tahun pajak ada 48 titik uji. Tiga kategori TER dengan 125 lapisan memberi 122 batas, atau 366 titik uji. Untuk sampel operasional, bila $n$ kasus lulus tanpa satu pun kegagalan, batas atas eksak laju kegagalan pada keyakinan $1-\alpha$ adalah $1 - \alpha^{1/n}$. Klaim laju kegagalan di bawah 0,1\% dengan keyakinan 95\% membutuhkan $n \ge 2.995$.

\begin{table}[t]
\centering
\caption{Ruang uji dan cakupannya}
\label{tab:cakupan}
\scriptsize
\begin{tabularx}{\columnwidth}{@{}L P{1.5cm} P{0.8cm} P{0.9cm}@{}}
\toprule
\textbf{Kriteria} & \textbf{Sel} & \textbf{Teruji} & \textbf{Acuan} \\
\midrule
Batas lapisan Pasal 17 & $4 \times 3 \times 4 = 48$ & \isian{} & \isian{} \\
Batas lapisan TER & $122 \times 3 = 366$ & \isian{} & \isian{} \\
Kombinasi 2-arah berbatasan & \isian{generator} & \isian{} & \isian{} \\
Kombinasi 3-arah berbatasan & \isian{generator} & \isian{} & \isian{} \\
Sampel operasional acak & $\ge 2.995$ & \isian{} & \isian{} \\
\bottomrule
\end{tabularx}
\end{table}

\definisikan{Generator kasus untuk ketiga kriteria, sumber distribusi sampel operasional, dan satuan titik uji Pasal 17 setelah pembulatan PKP ke ribuan.}

\subsection{Korpus}
\label{subsec:korpus}

\subsubsection{Peristiwa perubahan nyata}
Korpus memuat seluruh peristiwa pada 2023--2026 yang memengaruhi PPh 21 pegawai tetap, termasuk perubahan parameter tahunan seperti batas upah iuran pensiun. Tabel~\ref{tab:korpus} memuat kandidatnya. Hasil dilaporkan per jenis perubahan.

\begin{table}[t]
\centering
\caption{Kandidat korpus peristiwa perubahan}
\label{tab:korpus}
\scriptsize
\begin{tabularx}{\columnwidth}{@{}P{0.35cm} L P{1.25cm} P{1.6cm} P{0.45cm}@{}}
\toprule
\textbf{No.} & \textbf{Dokumen} & \textbf{Berlaku} & \textbf{Jenis} & $|\Delta^{\emas}|$ \\
\midrule
1 & Batas upah JP 2023 (surat BPJS TK) & 1 Mar 2023 & parameter & \isian{} \\
2 & PMK 66/2023 & 1 Jul 2023 & klasifikasi & \isian{} \\
3 & PP 58/2023 dan PMK 168/2023 & 1 Jan 2024 & metode, tabel tarif, pencabutan & \isian{} \\
4 & Batas upah JP 2024 & 1 Mar 2024 & parameter & \isian{} \\
5 & PMK 10/2025 (terbit 4 Feb 2025) & 1 Jan 2025 & fasilitas baru, berlaku surut & \isian{} \\
6 & Batas upah JP 2025 & 1 Mar 2025 & parameter & \isian{} \\
7 & PMK 72/2025 & Okt 2025 & kondisi atau cakupan & \isian{} \\
8 & PMK 105/2025 & 1 Jan 2026 & fasilitas baru & \isian{} \\
9 & Batas upah JP 2026 & 1 Mar 2026 & parameter & \isian{} \\
\midrule
K1 & PP 7/2025 dan PP 36/2025 (keringanan iuran JKK padat karya) & Feb 2025 & parameter, kondisi & \isian{} \\
\midrule
P1 & PMK 81/2024, PER-11/PJ/2025 (administrasi Coretax) & -- & pengecoh & 0 \\
P2 & PMK 6/2026 (DTP peserta pemagangan) & -- & pengecoh & 0 \\
P3 & PP 20/2026 (PPh final UMKM) & -- & pengecoh & 0 \\
\bottomrule
\end{tabularx}
\end{table}

\status{Baris 1--9 dan P1--P3 berasal dari catatan revisi. Baris K1 adalah kandidat tambahan dari \texttt{dataset/01\_regulasi/README.md}. Keringanan iuran JKK menurunkan premi JKK yang dibayar pemberi kerja, dan premi itu termasuk bruto, sehingga PPh 21 pegawai di perusahaan tersebut ikut berubah. Seluruh baris masih harus diverifikasi saat korpus disusun (Tahap 6).}

\subsubsection{Dokumen pengecoh}
Korpus juga memuat dokumen pengecoh, yaitu regulasi yang terbit pada periode yang sama tetapi tidak mengubah perhitungan PPh 21 pegawai tetap. Pengecoh sulit adalah dokumen yang menyebut PPh 21 tetapi hanya mengatur hal lain, misalnya pegawai tidak tetap, peserta pemagangan, atau administrasi pelaporan. Jawaban yang benar untuk pengecoh adalah pembaruan kosong.

\subsubsection{Dokumen kontrafaktual}
Generator membuat dokumen kontrafaktual dengan mengubah nilai, tanggal, atau ambang di dokumen nyata secara konsisten, lalu menurunkan $\Delta^{\emas}$ kontrafaktual secara terprogram dari $\Delta^{\emas}$ asli. Dokumen ini memeriksa apakah LLM membaca dokumen atau mengingat regulasi dari data latihnya. Varian kontrafaktual dihitung satu kelompok dengan dokumen sumbernya, sehingga menambah presisi, bukan jumlah sampel independen.

\subsubsection{Perubahan tak terlihat}
Perubahan struktural sintetis ditulis orang lain (Bagian~\ref{subsec:hindsight}), misalnya kategori TER keempat, jenis pengurang baru, tunjangan berplafon, aturan pembulatan baru, atau tarif harian.

\subsubsection{Kebijakan perusahaan}
Inti korpus kebijakan diambil dari dokumen nyata yang terbuka untuk umum, seperti perjanjian kerja bersama atau peraturan perusahaan yang dipublikasikan. Kebijakan buatan melengkapi jenis konflik yang tidak muncul di dokumen nyata. Konflik buatan ditulis orang lain dari taksonomi konflik sebelum detektor dibekukan, dan kedua kelompok dilaporkan terpisah. Setiap dokumen kebijakan juga menjadi masukan akuisisi (E2) dan diuji ekspresivitasnya (E1), dengan $\Delta^{\emas}$ yang disusun dan diperiksa dengan cara yang sama seperti dokumen regulasi. Perubahan kebijakan tidak disisipkan ke rangkaian \en{regulation replay}.

\status{Korpus kebijakan nyata saat ini kosong. Perusahaan X adalah rekonstruksi dari \en{spreadsheet} kantor, dan katalog KP-01 sampai KP-20 ditulis peneliti. Kandidat sumber dokumen nyata adalah koleksi PKB Indonesia di situs WageIndicator (sekitar 238 PKB, sekitar 37 di antaranya berlaku dalam 2023--2026, sebagian berpasangan versi lama dan baru). Lisensi redistribusinya belum diperiksa, dan sebagian PKB tidak memuat klausul PPh 21 sama sekali.}

\subsection{Mode evaluasi dan hasil per dokumen}
\label{subsec:mode}

Eksperimen akuisisi berjalan dalam dua mode.
\begin{itemize}[leftmargin=*,nosep]
  \item \textbf{Mode terisolasi (RQ2).} $\hat{\Delta}_n$ diterapkan pada $\K^{\emas}_{n-1}$.
  \item \textbf{Mode berantai (RQ3).} $\hat{\Delta}_n$ diterapkan pada $\K^{\mathrm{sis}}_{n-1}$, yaitu hasil pembaruan sistem sendiri pada langkah sebelumnya.
\end{itemize}
Kedua mode berjalan tanpa peninjau manusia, sehingga kandidat terpilih langsung diterapkan. Hasilnya adalah batas bawah kinerja, dan peran peninjau diukur lewat upaya koreksi. Peneliti tidak bertindak sebagai peninjau karena mengetahui jawaban emas.

Setiap dokumen relevan berakhir pada satu dari tiga hasil, yaitu benar-diterapkan, salah-diterapkan (kesalahan senyap), atau \en{abstain}. $\hat{\Delta}$ benar bila A2, A3$_t$, dan A4 terpenuhi pada seluruh $\SE$, sesuai Definisi~\ref{def:adapt-op}. Dokumen relevan yang dinilai tidak relevan oleh sistem dihitung salah-diterapkan. Dokumen pengecoh tidak masuk SER dan CAR. Penilaian relevansi diukur terpisah lewat laju \en{false positive} pada pengecoh dan laju pembaruan terlewat pada dokumen relevan.

Pada mode berantai, \en{abstain} ditangani dengan dua varian yang didaftarkan sebelumnya. Pada varian pertama, \en{knowledge base} lama dipertahankan, sehingga peristiwa berikutnya mewarisi kekosongan itu. Pada varian kedua, peninjau sempurna menyisipkan $\Delta^{\emas}_n$, sehingga perambatan kesalahan dapat diukur tanpa tercampur akibat \en{abstain}. Analisis konfirmatori memakai jalan pertama alur pada setiap dokumen sumber nyata. Jalan tambahan dan varian kontrafaktual dilaporkan sebagai strata terpisah.

\subsection{Pembanding dan ablasi}
\label{subsec:pembanding}

Tabel~\ref{tab:pembanding} memuat pembanding dan ablasi. B3 dan B4 memakai anggaran komputasi yang setara dengan usulan, atau dibandingkan secara eksplisit, sehingga keunggulan usulan tidak dapat dijelaskan oleh jumlah sampel saja. AB5 menjawab dengan angka pertanyaan mengapa \en{knowledge graph} diperlukan, dengan membandingkan validasi skema biasa dengan \en{shape} graf, dan pemindaian dependensi datar dengan \en{query} graf.

\begin{table}[t]
\centering
\caption{Pembanding dan ablasi}
\label{tab:pembanding}
\scriptsize
\begin{tabularx}{\columnwidth}{@{}P{0.55cm} L P{0.9cm}@{}}
\toprule
\textbf{ID} & \textbf{Kondisi} & \textbf{RQ} \\
\midrule
\multicolumn{3}{@{}l}{\textit{Pembanding}} \\
B1 & Praktik saat ini, diukur dari riwayat \en{commit} kalkulator PPh 21 \en{open source} di sekitar pergantian ke TER: waktu sampai diperbarui, baris yang berubah, dan \en{bug} regresi & motivasi, RQ3 \\
B2 & LLM menghitung PPh 21 langsung dari dokumen pada $\pm$100 kasus operasional dengan satu model & motivasi \\
B3 & Satu kandidat dengan gerbang G1--G2 saja (\en{pass@1}) & RQ2, RQ3 \\
B4 & $m(1+\bar{r})$ kandidat, sama dengan jumlah panggilan usulan termasuk putaran perbaikan, dipilih dengan suara terbanyak tanpa gerbang & RQ2 \\
B5 & Analisis dampak dengan tiga cara: menandai semua kebijakan yang menyentuh fakta yang berubah, menelusuri dependensi tanpa eksekusi, dan meminta penilaian LLM & RQ4 \\
\midrule
\multicolumn{3}{@{}l}{\textit{Ablasi (satu model)}} \\
AB1 & Formalisasi dari nol pada 2--3 keadaan hukum terpilih (sebelum TER, sesudah TER, masa DTP) dengan gerbang dan model yang sama; dinilai dengan A2 dan hilangnya riwayat & RQ2, RQ3 \\
AB2 & Tanpa konteks \en{knowledge base} & RQ2 \\
AB3 & Tanpa G5, tanpa G4, atau tanpa putaran perbaikan & RQ2 \\
AB4 & Pemilihan pertama lolos, konsensus, atau pilihan sempurna & RQ2 \\
AB5 & Skema biasa dibanding \en{shape} graf; pemindaian datar dibanding \en{query} graf & RQ2, RQ4 \\
\bottomrule
\end{tabularx}
\end{table}

\terbuka{B2 bertentangan dengan keputusan awal proyek untuk tidak memakai LLM sebagai pembanding. Usulan \en{skeleton}: B2 tetap disertakan dalam skala kecil sebagai pembanding motivasi, karena penelaah akan menanyakannya.}

\subsection{Hipotesis dan metrik utama}
\label{subsec:hipotesis}

Tabel~\ref{tab:lacak} menghubungkan setiap RQ dengan hipotesis, metrik utama, cara analisis, dan eksperimennya.

\begin{table*}[t]
\centering
\caption{Keterlacakan RQ, hipotesis, dan metrik utama}
\label{tab:lacak}
\scriptsize
\begin{tabularx}{\textwidth}{@{}P{0.9cm} L P{4.0cm} P{3.6cm} P{1.2cm}@{}}
\toprule
\textbf{RQ} & \textbf{Hipotesis} & \textbf{Metrik utama} & \textbf{Analisis} & \textbf{Eksperimen} \\
\midrule
-- & Invarian A1 dan A3$_\omega$ tanpa pelanggaran selama evaluasi & jumlah pelanggaran & pemeriksaan & semua \\
RQ1 (K1) & H1: laju adaptivitas representasional pada perubahan tak terlihat $\geq \theta_1$ & RA; dilaporkan juga untuk perubahan nyata dan korpus kebijakan & proporsi dengan CI eksak Clopper--Pearson & E1 \\
RQ2 (K2) & H2a: batas atas laju kesalahan senyap usulan $\leq \theta_2$. H2b: CAR usulan lebih tinggi daripada B3 dan B4, dan SER usulan tidak lebih buruk dari keduanya dengan margin $\delta_S$ & SER dan CAR atas dokumen sumber nyata, per lapisan & CI eksak; selisih risiko berpasangan dengan CI Newcombe; non-inferioritas SER & E2 \\
RQ3 (K3) & H3a: laju adaptivitas operasional berantai $\geq \theta_3$. H3b: kesalahan terpicu tidak melebihi batas derau pembangkitan & OAR dan kesalahan terpicu $I$ & CI eksak; deskriptif atas $R$ pengulangan, dibanding batas derau & E3 \\
RQ4 (K1) & H4a: \en{recall} $\geq \theta_4$. H4b: $F_2$ usulan lebih tinggi daripada setiap pembanding B5 & \en{Recall} dan $F_2$ & \en{bootstrap} klaster per peristiwa; batas Clopper--Pearson tingkat pasangan & E4 \\
\bottomrule
\end{tabularx}
\end{table*}

Metrik utama didefinisikan sebagai berikut.
\begin{itemize}[leftmargin=*,nosep]
  \item \textbf{RA (RQ1).} Untuk setiap perubahan nyata, perubahan tak terlihat, dan kebijakan dalam korpus kebijakan, RA mencatat apakah ada $\Delta$ dalam bahasa representasi yang memenuhi A1--A4. Ketiga kelompok dilaporkan terpisah, dan hanya perubahan tak terlihat yang menguji H1.
  \item \textbf{SER dan CAR (RQ2).} Keduanya dihitung atas himpunan $D$ dokumen sumber nyata yang relevan, yaitu dokumen regulasi dan dokumen kebijakan nyata, dengan satu hasil per dokumen, dan dilaporkan per lapisan. Varian kontrafaktual, pengecoh, dan kebijakan buatan tidak termasuk. Dengan $\mathbf{1}[\cdot]$ fungsi indikator,
  \begin{align}
  \mathrm{SER} &= \tfrac{1}{|D|}\textstyle\sum_{d \in D} \mathbf{1}[d \text{ salah-diterapkan}], \label{eq:ser}\\
  \mathrm{CAR} &= \tfrac{1}{|D|}\textstyle\sum_{d \in D} \mathbf{1}[d \text{ benar-diterapkan}], \nonumber
  \end{align}
  dan laju \en{abstain} $\mathrm{ABS} = 1 - \mathrm{SER} - \mathrm{CAR}$. SER mengukur keamanan dan CAR mengukur kegunaan. Keduanya dilaporkan bersama, karena sistem yang selalu \en{abstain} punya SER nol tetapi tidak berguna. Batas atas eksak untuk $f$ kesalahan dari $|D|$ dokumen adalah kuantil Beta $(1-\alpha;\ f+1,\ |D|-f)$, yang untuk $f=0$ menjadi $1-\alpha^{1/|D|}$.
  \item \textbf{OAR (RQ3).} Proporsi peristiwa yang tidak memperkenalkan kesalahan baru dalam mode berantai. A2, A3$_t$, dan A4 dinilai pada $\SE_{\leq n}$, yaitu gabungan kasus seluruh keadaan sampai peristiwa $n$ termasuk masa lampau, di luar tupel $H_n$ yang sudah salah sebelumnya. Tanpa pengecualian $H_n$, satu kesalahan lama terhitung ulang di setiap peristiwa berikutnya. Keadaan kumulatif dilaporkan terpisah lewat \en{drift} akhir $D_N$.
  \item \textbf{Kesalahan terpicu (RQ3).} Untuk peristiwa $n$, misalkan $E^{\mathrm{ber}}_n$ dan $E^{\mathrm{iso}}_n$ adalah himpunan tupel salah pada mode berantai dan mode terisolasi, dan $H_n$ himpunan tupel yang sudah salah di $\K^{\mathrm{sis}}_{n-1}$. Kesalahan terpicu adalah
  \begin{equation}
  I_n = E^{\mathrm{ber}}_n \setminus \big(E^{\mathrm{iso}}_n \cup H_n\big).
  \label{eq:terpicu}
  \end{equation}
  $I_n$ dipecah menjadi dua komponen. Komponen mekanis diperoleh dengan menerapkan $\hat{\Delta}_n$ dari mode terisolasi pada $\K^{\mathrm{sis}}_{n-1}$, sehingga yang terukur hanya perambatan lewat basis yang salah. Komponen kontekstual diperoleh dengan membangkitkan ulang $\hat{\Delta}_n$ memakai $\K^{\mathrm{sis}}_{n-1}$ sebagai konteks. Pembangkitan ulang memuat derau acak, sehingga batas derau diperkirakan dari selisih dua jalan terisolasi yang independen. H3b diuji terhadap batas itu, bukan terhadap nol.
  \item \textbf{\en{Recall} dan $F_2$ (RQ4).} Keduanya dihitung atas pasangan kebijakan dan peristiwa, dengan $F_2 = 5PR/(4P+R)$ untuk \en{precision} $P$ dan \en{recall} $R$. \en{Recall} diutamakan, karena kebijakan terdampak yang terlewat lebih berbahaya daripada peringatan palsu. Bila semua $n_+$ pasangan positif terdeteksi, batas bawah eksak \en{recall} adalah $\alpha^{1/n_+}$, misalnya $0{,}905$ untuk $n_+ = 30$. Pembanding yang menandai semua kebijakan menghasilkan $F_2 = 5\pi/(4\pi+1)$ untuk prevalensi $\pi$, sehingga $F_2$ baru bermakna di atas angka itu.
\end{itemize}

Ambang harus dapat dicapai oleh jumlah sampel yang direncanakan. Dengan sepuluh dokumen uji tanpa satu pun kesalahan, batas atas eksak SER pada keyakinan 95\% hanya $1-0{,}05^{1/10} \approx 0{,}26$. Klaim SER $\leq 5\%$ baru mungkin dengan minimal 59 dokumen tanpa kesalahan. Untuk setiap hipotesis, pra-registrasi mencantumkan batas terbaik yang dapat dicapai oleh jumlah sampel yang direncanakan, lalu menetapkan $\theta$ pada atau di atas batas itu. Bila tidak ada ambang yang bermakna, hipotesis diganti dengan estimasi dan CI saja. Klaim keamanan lalu dinyatakan apa adanya, misalnya SER paling tinggi 26\% dengan keyakinan 95\% pada sepuluh dokumen, dan diperketat lewat peristiwa prospektif.

\definisikan{Ambang $\theta_1$--$\theta_4$ dan margin $\delta_S$ beserta alasannya, dikunci sebelum eksperimen.}

\subsection{Metrik diagnostik}
\label{subsec:diagnostik}

Metrik di Tabel~\ref{tab:diagnostik} menjelaskan mengapa metrik utama bernilai seperti itu. Metrik ini eksploratori dan tidak dipakai untuk menguji hipotesis. Metrik \en{pass@k} dan \en{n@k} dipinjam dari evaluasi kode buatan LLM~\cite{chen2021evaluating,li2022alphacode}, dan kurva risiko terhadap cakupan dari prediksi selektif~\cite{geifman2017selective}.

\begin{table*}[t]
\centering
\caption{Metrik diagnostik (eksploratori)}
\label{tab:diagnostik}
\scriptsize
\begin{tabularx}{\textwidth}{@{}P{1.2cm} P{3.3cm} L P{3.6cm}@{}}
\toprule
\textbf{RQ} & \textbf{Metrik} & \textbf{Definisi} & \textbf{Pertanyaan yang dijawab} \\
\midrule
Instrumen & Selisih per hasil adjudikasi; batas atas laju kegagalan & semua selisih $\K^{\emas}$ terhadap GT1--GT5; $p_U = 1-\alpha^{1/n}$ untuk $n$ kasus operasional tanpa kegagalan & seberapa dapat dipercaya \en{ground truth} \\
Instrumen & \en{Mutation score} global dan terfokus-$\Delta$ & mutan terdeteksi dibagi mutan tak setara & apakah $\SE$ cukup tajam \\
RQ1 & Akurasi resolusi konflik & pasangan aturan dengan pemenang hukum yang diketahui per asas, dinilai benar atau salah & apakah resolusi konflik bekerja \\
RQ1 & Kelengkapan jejak & fakta keluaran yang jejaknya sampai ke sumber hukum, dibagi seluruh fakta keluaran & kelengkapan fasilitas penjelasan \\
\midrule
RQ2 & \en{pass@1}, \en{pass@m}, \en{1@m}, efisiensi pemilihan & pada kandidat sebelum perbaikan: $\text{pass@}j = 1-\binom{m-c}{j}/\binom{m}{j}$; $\eta = (\text{1@}m - \text{pass@}1)/(\text{pass@}m - \text{pass@}1)$ & seberapa baik gerbang memilih dibanding pemilihan sempurna \\
RQ2 & Laju diterapkan tetapi tidak lengkap & $\hat{\Delta}$ diterapkan yang tidak memuat minimal satu unsur $\Delta^{\emas}$ & apakah kelalaian adalah kesalahan dominan \\
RQ2 & Matriks konfusi gerbang & per gerbang: kebocoran $=\mathrm{FN}/(\mathrm{TP}+\mathrm{FN})$, penolakan keliru $=\mathrm{FP}/(\mathrm{FP}+\mathrm{TN})$, dan MCC; kandidat salah sebagai kelas positif & gerbang mana yang efektif \\
RQ2 & Kesepakatan perilaku dan kebenaran per unsur & proporsi tupel $\SE$ yang sama dengan $\K^{\emas}_n$; unsur salah bila menukarnya dengan unsur emas mengubah hasil $\SE$ & seberapa dekat pembaruan yang salah \\
RQ2 & Laju hafalan dan kesetiaan kontrafaktual & proporsi nilai yang diubah di mana $\hat{\Delta}$ memuat nilai asli, dan yang memuat nilai kontrafaktual & apakah LLM membaca dokumen atau mengingat \\
RQ2 & Kontaminasi & SER dan CAR untuk peristiwa sebelum dan sesudah \en{training cutoff} model & pengaruh data latih \\
RQ2 & Penilaian relevansi & laju \en{false positive} pada pengecoh dan laju pembaruan terlewat pada dokumen relevan & perubahan tanpa alasan dan perubahan terlewat \\
RQ2 & Strata kontrafaktual dan kebijakan buatan & SER dan CAR pada kedua strata, CI \en{bootstrap} klaster per dokumen sumber & ketahanan terhadap nilai dan jenis kebijakan yang tidak dikenal LLM \\
RQ2 & Ketepatan kutipan & unsur yang kutipannya ada dan mendukung unsur itu & dasar hukum karangan \\
RQ2 & Upaya koreksi & operasi sunting atribut dari $\hat{\Delta}$ ke $\Delta^{\emas}$ setelah penyelarasan unsur, hanya pada unsur yang salah secara perilaku; penghematan $= 1 - \text{sunting}/\text{operasi menyusun } \Delta^{\emas} \text{ dari kosong}$ & pekerjaan manusia yang dihemat \\
RQ2 & Kurva biaya dan akurasi & CAR terhadap biaya per peristiwa & kelayakan praktis \\
\midrule
RQ3 & Rasio perambatan & $P(\hat{\Delta}_n \text{ salah} \mid \K^{\mathrm{sis}}_{n-1} \text{ salah}) / P(\hat{\Delta}_n \text{ salah} \mid \K^{\emas}_{n-1})$ & apakah basis yang salah memperbesar kesalahan berikutnya \\
RQ3 & Formalisasi dari nol (AB1) & A2 pada masa terpilih dan hilangnya riwayat, dibanding usulan & apa yang hilang bila \en{knowledge base} dibangun ulang \\
RQ3 & Persistensi, pemulihan, penguncian & jumlah peristiwa yang dilalui kesalahan sebelum hilang; kesalahan bawaan yang perbaikannya ditolak G4 & apakah kesalahan menetap atau terawetkan gerbang \\
RQ3 & \en{Drift} akhir dan anomali & $D_N$: proporsi pegawai-bulan berbobot yang PPh 21-nya salah setelah seluruh rangkaian, beserta rupiahnya; aturan redundan, berkonflik, tercakup, mati, rujukan menggantung & akibat finansial akhir dan mutu \en{knowledge base} \\
RQ3 & Penerapan langsung & peristiwa yang diterapkan tanpa suntingan manusia, dengan paling banyak $k$ suntingan, dan waktu dari dokumen tersedia sampai \en{knowledge base} benar & bukti langsung bahwa perubahan regulasi dapat diterapkan \\
\midrule
RQ4 & \en{Precision}, MCC, AUPRC, $\kappa$ & dengan masukan $\Delta^{\emas}$ (detektor saja) dan $\hat{\Delta}$ (alur penuh); AUPRC untuk pembanding yang menghasilkan skor; Cohen's $\kappa$ untuk label jenis (c) & sumber kesalahan analisis dampak dan keandalan label \\
\bottomrule
\end{tabularx}
\end{table*}

\subsection{Analisis statistik}
\label{subsec:statistik}

Jumlah peristiwa nyata kecil. Karena itu penelitian ini tidak memakai uji yang membutuhkan sampel besar, dan melaporkan hasil sebagai estimasi dengan selang kepercayaan eksak beserta efek minimum yang dapat dideteksi (Tabel~\ref{tab:statistik}).

\begin{table*}[t]
\centering
\caption{Rencana analisis statistik}
\label{tab:statistik}
\scriptsize
\begin{tabularx}{\textwidth}{@{}P{3.3cm} P{3.6cm} L P{2.7cm}@{}}
\toprule
\textbf{Jenis data} & \textbf{Contoh} & \textbf{Penaksir atau uji} & \textbf{Ukuran efek} \\
\midrule
Pemeriksaan nol-cacat & invarian A1 dan A3$_\omega$; selisih $\K^{\emas}$ & hitungan dan batas atas eksak & -- \\
Proporsi per kondisi & RA, SER, CAR, OAR, \en{recall} & CI eksak Clopper--Pearson~\cite{clopper1934confidence} & proporsi \\
Biner berpasangan antarkondisi & benar-diterapkan, usulan dibanding B3 dan B4 & selisih risiko berpasangan dengan CI Newcombe~\cite{newcombe1998improved}; McNemar eksak sebagai pelengkap & selisih risiko; \en{odds ratio} pasangan sumbang \\
Non-inferioritas & SER usulan dibanding B3 dan B4 & batas atas CI selisih SER berpasangan $< \delta_S$ & selisih risiko \\
Kontinu berpasangan & kesepakatan perilaku, upaya koreksi & Wilcoxon bertanda & \en{rank-biserial} $r$ \\
Antarpengulangan rangkaian & OAR dan $I$ atas $R$ pengulangan & median dan rentang dibanding batas derau; uji permutasi bila perlu & -- \\
Proporsi berklaster & \en{recall} dan $F_2$ RQ4; strata kontrafaktual & \en{bootstrap} klaster per peristiwa atau per dokumen sumber & selisih $F_2$ \\
Kesepakatan & label RQ4 dan $\Delta^{\emas}$ & Cohen's $\kappa$~\cite{cohen1960coefficient}; Krippendorff's $\alpha$~\cite{hayes2007answering} & $\kappa$, $\alpha$ \\
\bottomrule
\end{tabularx}
\end{table*}

Unit analisisnya adalah dokumen sumber nyata untuk RQ2, dan peristiwa untuk RQ3 dan RQ4. Jalan LLM, kandidat, dan varian kontrafaktual bersarang di dalam dokumen sumbernya dan tidak dihitung sebagai sampel independen. Hasil RQ3 bersyarat pada satu rangkaian nyata. $R$ pengulangan hanya menangkap keragaman pembangkitan LLM, sehingga hasilnya dilaporkan secara deskriptif. Koreksi Holm~\cite{holm1979simple} diterapkan per RQ dan hanya pada hipotesis komparatif. Pemeriksaan nol-cacat dan invarian tidak termasuk. Hasil dilaporkan per jenis perubahan, dan jenis mudah tidak digabung dengan jenis struktural.

\definisikan{Nilai $m$ (usulan 5), $R$ (usulan 3), parameter pembangkitan, jumlah dokumen per bagian data, dan efek minimum yang dapat dideteksi.}

\subsection{Prosedur eksperimen}
\label{subsec:prosedur}

\begin{itemize}[leftmargin=*,nosep]
  \item \textbf{E0, validasi instrumen.} Bagian~\ref{subsec:instrumen}. E0 menjadi prasyarat semua eksperimen lain.
  \item \textbf{E1, adaptivitas representasi (RQ1).} RA pada perubahan nyata, perubahan tak terlihat, dan korpus kebijakan; akurasi resolusi konflik; dan pemetaan riwayat \en{engine} ke regulasi pemicunya.
  \item \textbf{E2, akuisisi terisolasi (RQ2).} Setiap dokumen uji, baik regulasi maupun kebijakan, dijalankan pada kondisi usulan, B3, B4, dan AB2--AB5. AB1 dijalankan pada keadaan hukum terpilih. Model utama dipakai untuk semua kondisi, dan model kedua hanya untuk usulan dan B3. Satu jalan terisolasi tambahan dipakai untuk batas derau.
  \item \textbf{E3, \en{regulation replay} (RQ3).} Mode berantai dari \en{knowledge base} awal 2023 sampai peristiwa terakhir, dalam dua varian \en{abstain} dan $R$ pengulangan. Kesalahan didekomposisi menjadi bawaan, baru, terpicu mekanis, dan terpicu kontekstual. B1 menjadi pembanding praktik, dan peristiwa prospektif ditambahkan begitu terbit.
  \item \textbf{E4, analisis dampak (RQ4).} Dua masukan, yaitu $\Delta^{\emas}$ dan $\hat{\Delta}$, dibandingkan dengan pembanding B5.
  \item \textbf{E5, motivasi.} B2 pada $\pm$100 kasus operasional.
\end{itemize}

\subsection{Biaya, kolaborator, dan reproduksibilitas}
\label{subsec:biaya}

Jumlah panggilan LLM dihitung per eksperimen. Misalkan $\bar{r}$ adalah rata-rata putaran perbaikan dan $D'$ gabungan $D$, dokumen kontrafaktual, dan pengecoh. E2 membutuhkan $|D'| \times m(1+\bar{r})$ panggilan untuk usulan, $|D'|$ untuk B3, dan $|D'| \times m(1+\bar{r})$ untuk B4. Setiap varian AB2--AB5 membutuhkan $|D| \times m(1+\bar{r})$. AB1 membutuhkan jumlah keadaan hukum terpilih dikali $m(1+\bar{r})$ dengan konteks seluruh dokumen yang berlaku, dan inilah panggilan termahal. E3 membutuhkan $2 \times R \times N \times m(1+\bar{r})$ untuk dua varian \en{abstain}, ditambah $|D| \times m(1+\bar{r})$ untuk batas derau. E5 membutuhkan sekitar 100 panggilan. Total panggilan dikalikan rata-rata token dan harga model, lalu dicantumkan dalam rupiah di pra-registrasi. Biaya ditekan dengan memotong dokumen per pasal, memakai \en{prompt caching}, membangkitkan kontrafaktual secara otomatis, dan menjalankan ablasi pada satu model.

Beberapa langkah membutuhkan orang selain peneliti, yaitu penulis perubahan tak terlihat dan konflik buatan, penulis contoh perhitungan kebijakan (GT5), pembaca kedua $\Delta^{\emas}$, pengadjudikasi kedua, pelabel kedua RQ4, dan praktisi pajak. Peran-peran ini digabung ke paling banyak dua kolaborator bernama, dan jam kerjanya dicatat. Untuk setiap peran yang tidak terisi, pra-registrasi mencantumkan apa yang dihapus dan bagaimana hal itu dilaporkan sebagai keterbatasan. Contohnya, tanpa penulis perubahan tak terlihat, RQ1 hanya melaporkan perubahan nyata, dan klaimnya turun menjadi kemampuan menyatakan perubahan yang sudah diketahui.

Versi perangkat lunak dan model serta tanggal akses dicatat, dan seluruh \en{prompt} serta keluaran LLM diarsipkan, sehingga semua analisis dapat diulang tanpa memanggil LLM lagi. Materi tersembunyi, yaitu perubahan tak terlihat, konflik buatan, GT5, subset \en{holdout}, dan nilai harapan $\SE$, disimpan di repositori privat selama penelitian. Hash SHA-256 berkasnya dicatat di repositori publik sebagai bukti waktu.

\terbuka{Kolaborator, paling banyak dua orang, untuk peran-peran di atas.}

% ============================================================================
\section{Hasil}
\label{sec:hasil}
% ============================================================================

\catatan{Belum ada hasil. Bagian ini hanya kerangka urutan pelaporan, tabel kosong, dan gambar pengganti. Isi setelah E0 lulus dan E1--E4 dijalankan sesuai pra-registrasi. Hipotesis yang terbantah dilaporkan sebagai terbantah, dan analisis di luar pra-registrasi diberi label eksploratori.}

Urutan pelaporan mengikuti Tabel~\ref{tab:lacak}. Validasi instrumen dilaporkan lebih dulu. Setiap RQ dibuka dengan metrik utama dan putusan atas hipotesisnya, lalu metrik diagnostik.

\subsection{Validasi instrumen}
\label{subsec:hasil-e0}
\isian{Selisih $\K^{\emas}$ terhadap GT1--GT5 menurut hasil adjudikasi; matriks acuan; Tabel~\ref{tab:cakupan} yang sudah terisi; \en{mutation score} global dan terfokus-$\Delta$; Krippendorff's $\alpha$ untuk $\Delta^{\emas}$; jumlah koreksi emas.}

\subsection{Adaptivitas representasi (RQ1)}
\label{subsec:hasil-rq1}
\isian{RA per kelompok dengan CI 95\% dan putusan H1; perubahan dan kebijakan yang gagal beserta alasannya; riwayat \en{engine} yang dipetakan ke regulasi pemicunya; akurasi resolusi konflik.}

\begin{table}[t]
\centering
\caption{Adaptivitas representasional (E1)}
\label{tab:hasil-rq1}
\scriptsize
\begin{tabularx}{\columnwidth}{@{}L P{0.5cm} P{1.0cm} P{1.9cm}@{}}
\toprule
\textbf{Kelompok} & $n$ & \textbf{Berhasil} & \textbf{RA (CI 95\%)} \\
\midrule
Perubahan nyata 2023--2026 & \isian{} & \isian{} & \isian{} \\
Perubahan tak terlihat (H1) & \isian{} & \isian{} & \isian{} \\
Kebijakan perusahaan nyata & \isian{} & \isian{} & \isian{} \\
Kebijakan buatan & \isian{} & \isian{} & \isian{} \\
\bottomrule
\end{tabularx}
\end{table}

\subsection{Akuisisi satu langkah (RQ2)}
\label{subsec:hasil-rq2}
\isian{SER, CAR, dan ABS per kondisi, per lapisan, dan per jenis perubahan dengan CI; selisih CAR dan non-inferioritas SER terhadap B3 dan B4 beserta putusan H2a dan H2b; penilaian relevansi; strata kontrafaktual; \en{pass@1}, \en{1@m}, dan \en{pass@m}; metrik diagnostik.}

\begin{table}[t]
\centering
\caption{Keamanan dan kegunaan akuisisi per kondisi (E2)}
\label{tab:hasil-rq2}
\scriptsize
\begin{tabularx}{\columnwidth}{@{}L P{1.45cm} P{1.45cm} P{1.1cm}@{}}
\toprule
\textbf{Kondisi} & \textbf{SER (CI)} & \textbf{CAR (CI)} & \textbf{ABS} \\
\midrule
Usulan, lapisan regulasi & \isian{} & \isian{} & \isian{} \\
Usulan, lapisan perusahaan & \isian{} & \isian{} & \isian{} \\
B3 (\en{pass@1}, G1--G2) & \isian{} & \isian{} & \isian{} \\
B4 (suara terbanyak) & \isian{} & \isian{} & \isian{} \\
AB2--AB5 & \isian{} & \isian{} & \isian{} \\
\bottomrule
\end{tabularx}
\end{table}

\begin{figure}[t]
\centering
\gambar{Setiap kondisi digambar sebagai titik dengan CI. Sumbu-x adalah cakupan, yaitu CAR $+$ SER. Sumbu-y adalah risiko selektif, yaitu SER$/$(CAR $+$ SER), atau laju kesalahan di antara pembaruan yang diterapkan.}
\caption{Keamanan dan kegunaan akuisisi per kondisi.}
\label{fig:rq2}
\end{figure}

\subsection{Adaptivitas operasional berantai (RQ3)}
\label{subsec:hasil-rq3}
\isian{OAR dan putusan H3a; kesalahan bawaan, baru, dan terpicu per peristiwa beserta putusan H3b; rasio perambatan; jumlah peristiwa yang diterapkan tanpa suntingan manusia dan waktunya; \en{drift} akhir dalam rupiah; anomali \en{knowledge base}; perbandingan dengan B1.}

\begin{figure*}[t]
\centering
\gambarlebar{Satu batang bertumpuk per peristiwa, diurutkan kronologis di sumbu-x dan diberi label jenis perubahan. Batang menunjukkan kesalahan bawaan, baru, dan terpicu. Penanda menunjukkan status A2, A3$_t$, dan A4, dan garis menunjukkan rentang antarpengulangan.}
\caption{Hasil \en{regulation replay}.}
\label{fig:replay}
\end{figure*}

\subsection{Analisis dampak lintas lapis (RQ4)}
\label{subsec:hasil-rq4}
\isian{\en{Recall} dan $F_2$ dengan CI pada kedua masukan ($\Delta^{\emas}$ dan $\hat{\Delta}$) beserta putusan H4a dan H4b; perbandingan dengan B5, termasuk garis dasar \emph{tandai semua}; Cohen's $\kappa$; satu contoh penelusuran dari graf.}

\subsection{Ablasi}
\label{subsec:hasil-ablasi}
\isian{Perubahan SER dan CAR untuk setiap ablasi. AB5 menjawab dengan angka mengapa \en{knowledge graph} diperlukan.}

% ============================================================================
\section{Pembahasan}
\label{sec:pembahasan}
% ============================================================================

\catatan{Subbagian A--C menunggu hasil. Subbagian D dan E sudah dapat ditulis dari rancangan, dan akan disesuaikan dengan hasil.}

\subsection{Jawaban atas pertanyaan penelitian}
\isian{Satu paragraf per RQ, berisi metrik utama dengan CI dan putusan hipotesisnya.}

\subsection{Analisis kesalahan LLM}
Taksonomi kesalahan ditetapkan sebelum pelabelan, agar label tidak menyesuaikan diri dengan keluaran yang sudah terlihat. Jenisnya adalah nilai salah, masa berlaku salah, kondisi atau cakupan salah, pengecualian terlewat, pencabutan terlewat, lapisan atau sifat salah, kutipan karangan, dan perubahan pada pengecoh. \isian{Frekuensi per jenis, gerbang yang menangkap setiap jenis, dan contoh kesalahan yang lolos semua gerbang.}

\subsection{Implikasi}
\isian{Implikasi bagi pengembang sistem payroll dan bagi formalisasi peraturan yang sering berubah.}

\subsection{Ancaman terhadap validitas}
\label{subsec:ancaman}

\textbf{Validitas internal.} \en{Ground truth} dan \en{engine} disusun tim yang sama, sehingga kesalahan yang sama dapat muncul di keduanya. Mitigasinya adalah acuan pihak ketiga, subset \en{holdout} contoh resmi, pengadjudikasi kedua, dan uji \en{engine} yang terpisah dari \en{knowledge base}.

\textbf{Bias pandangan belakang.} Representasi dirancang setelah perubahan 2023--2026 diketahui. Mitigasinya adalah perubahan tak terlihat yang ditulis orang lain, peristiwa prospektif, dan pelaporan riwayat \en{engine}.

\textbf{Validitas konstruk.} Kecocokan dengan acuan tidak sama dengan kesesuaian hukum. Penelitian ini tidak melibatkan pakar hukum pajak, kecuali bila praktisi pajak dapat dilibatkan untuk memeriksa sampel. Upaya manusia diukur dengan proksi, yaitu jumlah operasi sunting.

\textbf{Validitas eksternal.} Penelitian ini mencakup satu yurisdiksi, satu jenis pajak, dan satu rangkaian perubahan. Hasilnya belum tentu berlaku untuk pajak lain atau negara lain.

\textbf{Validitas kesimpulan.} Jumlah peristiwa nyata kecil, LLM tidak deterministik, dan versi model berubah dari waktu ke waktu. Karena itu hasil dilaporkan sebagai estimasi dengan selang kepercayaan eksak, dan seluruh masukan serta keluaran LLM diarsipkan.

\textbf{Kontaminasi data.} Regulasi 2023--2026 kemungkinan ada di data latih LLM. Mitigasinya adalah dokumen kontrafaktual, perbandingan hasil sebelum dan sesudah \en{training cutoff} model, dan peristiwa prospektif.

\subsection{Keterbatasan}
\label{subsec:keterbatasan}

Penelitian ini terbatas pada pegawai tetap dengan satu pemberi kerja. Pegawai tidak tetap, pensiunan, pemberi kerja ganda, dan fasilitas khusus Ibu Kota Nusantara tidak dicakup, dan nilai natura diterima sebagai masukan. Jenis perubahan yang tidak terjadi pada 2023--2026, seperti perubahan PTKP, hanya terwakili lewat data pengembangan dan perubahan tak terlihat. \isian{Keterbatasan lain dari hasil, misalnya jenis perubahan yang selalu membuat sistem \en{abstain}, dan peran independen yang tidak terisi.}

% ============================================================================
\section{Kesimpulan dan penelitian lanjutan}
\label{sec:kesimpulan}
% ============================================================================

\isian{Jawaban RQ1--RQ4 dengan metrik utama dan CI, tanpa mengulang abstrak.}

\isian{Penelitian lanjutan yang diturunkan dari keterbatasan. Kandidatnya antara lain perluasan ke PPh 21 bukan pegawai dan PPh 26, pemberi kerja ganda, penambahan peristiwa prospektif, dan penerapan protokol \en{regulation replay} pada yurisdiksi lain.}

% ----------------------------------------------------------------------------
\section*{Pengungkapan penggunaan AI}
% ----------------------------------------------------------------------------
Penyusunan draf naskah ini dibantu asisten AI Claude (model Claude Opus 5.5, Anthropic) melalui Claude Code, untuk menyusun kalimat dari kerangka penelitian dan menyiapkan kode \LaTeX{}. Seluruh isi, klaim, angka, dan rujukan diperiksa dan menjadi tanggung jawab penulis. Riwayat penggunaan AI dicatat di \texttt{journal.md} repositori penelitian.

\catatan{Sesuaikan pernyataan ini dengan kebijakan penerbit tujuan sebelum naskah dikirim. IEEE meminta pengungkapan konten buatan AI di bagian ucapan terima kasih.}

% ----------------------------------------------------------------------------
\ifrancangan
\section*{Keputusan terbuka sebelum pra-registrasi}
\kotak{orange}{Daftar dari \en{skeleton}, Bagian X}{
\begin{enumerate}[leftmargin=*,nosep]
  \item Target publikasi. Bila IEEE Access, naskah final harus berbahasa Inggris.
  \item Apakah \en{preprint} boleh dikutip.
  \item Ukuran \en{lex specialis}.
  \item Model utama, model kedua, dan anggaran API.
  \item Apakah pembanding B2 (LLM menghitung langsung) disertakan.
  \item Apakah LLM boleh memakai alat ekstraksi tabel.
  \item Kolaborator, paling banyak dua, untuk perubahan tak terlihat, konflik buatan, contoh perhitungan kebijakan, pembaca kedua, dan pelabel kedua, beserta rencana cadangan bila tidak ada.
  \item Ambang hipotesis yang dapat dicapai oleh jumlah dokumen dan peristiwa yang tersedia.
  \item Pembaca kedua $\Delta^{\emas}$ dan praktisi pajak.
  \item Tambahan draf ini: lingkup \en{gross-up} rezim disetahunkan, sifat \emph{tafsir} sebagai sifat aturan keempat, dan aturan bila dua kelompok konsensus sama besar.
\end{enumerate}}
\fi


\bibliographystyle{IEEEtran}
\bibliography{../../latex/references,../../latex/regulasi,rujukan_tambahan}

\end{document}
